\documentclass{AnantDocumentation}

\title{Compre Notes for DIP}
\author{Pranav}

\begin{document}
\maketitle
\tableofcontents
\section{Convolution}

\subsection{Theory}

Convolution is a fundamental operation in linear image processing, describing how a spatial filter (kernel) modifies an image.  
It represents the weighted sum of pixel intensities within a neighborhood, where the weights are given by the filter coefficients.  

Formally, convolution models the response of a linear, shift-invariant (LSI) system when an image $f(x, y)$ passes through an impulse response $h(x, y)$.

\subsection{Mathematical Formulation}

For continuous images:
\begin{equation}
g(x, y) = (f * h)(x, y) = \int_{-\infty}^{\infty}\int_{-\infty}^{\infty} f(\xi, \eta)\, h(x - \xi, y - \eta) \, d\xi \, d\eta
\end{equation}

For discrete digital images:
\begin{equation}
g[m, n] = \sum_{s=-a}^{a}\sum_{t=-b}^{b} f[m - s, n - t] \, h[s, t]
\end{equation}
where $h[s, t]$ is the filter mask of size $(2a + 1) \times (2b + 1)$.

\paragraph{Properties:}
\begin{itemize}
\item \textbf{Commutative:} $f * h = h * f$
\item \textbf{Associative:} $(f * h_1) * h_2 = f * (h_1 * h_2)$
\item \textbf{Distributive:} $f * (h_1 + h_2) = f * h_1 + f * h_2$
\item \textbf{Shift Invariance:} Translating the input shifts the output by the same amount.
\end{itemize}

\paragraph{Intuitive Meaning:}
Each output pixel $g[m, n]$ is computed as a local weighted average of input pixels, centered at $(m, n)$, using the flipped kernel $h[-s, -t]$.

\subsection{Example 1: Smoothing by Convolution}

Let
\[
f =
\begin{bmatrix}
1 & 2 & 3 \\
4 & 5 & 6 \\
7 & 8 & 9
\end{bmatrix}, \quad
h = \frac{1}{9}
\begin{bmatrix}
1 & 1 & 1 \\
1 & 1 & 1 \\
1 & 1 & 1
\end{bmatrix}
\]

Then
\[
g(2,2) = \frac{1}{9} \sum_{i=-1}^{1}\sum_{j=-1}^{1} f(2-i,2-j)
= \frac{1}{9}(1+2+3+4+5+6+7+8+9) = 5
\]

Thus, each pixel becomes the average of its $3\times3$ neighborhood.  
This operation smooths the image and suppresses noise, corresponding to a low-pass filter in the frequency domain.

\subsection{Example 2: Edge Detection by Convolution}

Consider the Sobel horizontal edge detector:
\[
h =
\begin{bmatrix}
-1 & 0 & 1 \\
-2 & 0 & 2 \\
-1 & 0 & 1
\end{bmatrix}
\]
When applied to an image, the convolution emphasizes regions with large horizontal intensity gradients.

For a row
$f = [10,\, 20,\, 40]$,  
the convolution result (1D simplification) is:
\[
g = (-1)(10) + (0)(20) + (1)(40) = 30
\]
A large magnitude indicates a strong edge at that position.

\subsection{Summary}
Convolution is a mathematical model for spatial filtering that:
\begin{itemize}
\item Describes linear, shift-invariant image transformations.
\item Links directly to multiplication in the frequency domain via the Convolution Theorem.
\item Serves as the computational foundation for blurring, sharpening, and noise removal operations.
\end{itemize}
\section{Filtering}

\subsection{Theory}

Filtering refers to the process of modifying or enhancing an image by operating on its pixel values using a neighborhood-based operator.  
Filters can be classified broadly into:

\begin{itemize}
\item \textbf{Linear filters:} The output is a weighted sum of input pixels (e.g., averaging, Gaussian).
\item \textbf{Nonlinear filters:} The output is computed via nonlinear operations (e.g., median, max/min).
\end{itemize}

In image processing, filtering serves multiple purposes:
\begin{itemize}
\item Noise reduction (smoothing)
\item Edge enhancement (sharpening)
\item Feature extraction (edge/line detection)
\item Texture analysis
\item Preprocessing for segmentation or compression
\end{itemize}

Filtering is implemented either in the \textbf{spatial domain} (mask-based operations) or the \textbf{frequency domain} (multiplication with a frequency response).

\subsection{Mathematical Formulation}

\subsubsection{Linear Filtering}

A linear spatial filter computes each output pixel as:
\begin{equation}
g[m,n] = \sum_{s=-a}^{a} \sum_{t=-b}^{b} w[s,t] \, f[m-s, n-t]
\end{equation}
where $w[s,t]$ is the filter mask and $(2a+1)\times (2b+1)$ is its size.

This is equivalent to convolution:
\[
g = f * w
\]

Examples include:
\begin{itemize}
\item Averaging filter
\item Gaussian smoothing filter
\item Laplacian sharpening
\item Sobel/Prewitt gradient filters
\end{itemize}

\subsubsection{Nonlinear Filtering}

Nonlinear filters do not compute weighted sums. Instead, they use ranking or other nonlinear rules:
\[
g[m,n] = \text{nonlinear}(f \text{ in neighborhood})
\]

Examples:
\begin{itemize}
\item \textbf{Median filter:}
\[
g[m,n] = \text{median}\{f[i,j] \mid (i,j)\in \mathcal{N}_{m,n}\}
\]
\item Min/max filtering
\item Order-statistic filters
\end{itemize}

Nonlinear filters are especially useful for removing \textbf{salt-and-pepper noise}.

\subsection{Important Types of Filters}

\subsubsection{Smoothing (Low-Pass) Filters}

Used to reduce high-frequency noise.

\paragraph{Averaging Filter:}
\[
w = \frac{1}{mn} \mathbf{1}_{m \times n}
\]

\paragraph{Gaussian Low-Pass Filter:}
\[
w(s,t) = \frac{1}{2\pi\sigma^2} e^{-\frac{s^2+t^2}{2\sigma^2}}
\]

This filter preserves edges better than simple averaging.

\subsubsection{Sharpening (High-Pass) Filters}

Highlight high-frequency content such as edges.

\paragraph{Laplacian Filter:}
\[
\nabla^2 f = 
\begin{bmatrix}
0 & 1 & 0 \\
1 & -4 & 1 \\
0 & 1 & 0
\end{bmatrix} * f
\]

Sharpened image:
\[
g = f - \lambda (\nabla^2 f)
\]

\paragraph{Gradient Filters (Sobel):}
\[
G_x =
\begin{bmatrix}
-1 & 0 & 1\\
-2 & 0 & 2\\
-1 & 0 & 1
\end{bmatrix}, \qquad 
G_y =
\begin{bmatrix}
-1 & -2 & -1\\
0 & 0 & 0\\
1 & 2 & 1
\end{bmatrix}
\]

\subsubsection{Bandpass and Bandreject Filters}

Used to isolate or remove specific frequency ranges.  
Mathematically defined in the frequency domain as:
\[
G(u,v) = H(u,v)F(u,v)
\]

Where $H(u,v)$ selectively passes or rejects frequency bands.

\subsection{Example Problem 1: Gaussian Smoothing}

\subsubsection*{Problem}
Given the $3 \times 3$ Gaussian filter with $\sigma=1$:
\[
w = \frac{1}{16}
\begin{bmatrix}
1 & 2 & 1\\
2 & 4 & 2\\
1 & 2 & 1
\end{bmatrix}
\]
Apply the filter to the $3 \times 3$ image:
\[
f =
\begin{bmatrix}
10 & 20 & 30\\
20 & 40 & 60\\
30 & 60 & 90
\end{bmatrix}
\]
Compute the filtered output at the center pixel.

\subsubsection*{Solution}
\[
g(2,2) = \frac{1}{16} (1\cdot 10 + 2\cdot 20 + 1\cdot 30
+ 2\cdot 20 + 4\cdot 40 + 2\cdot 60
+ 1\cdot 30 + 2\cdot 60 + 1\cdot 90)
\]

Compute numerator:
\[
10 + 40 + 30 + 40 + 160 + 120 + 30 + 120 + 90 = 640
\]

Thus:
\[
g(2,2) = \frac{640}{16} = 40
\]

\subsubsection*{Interpretation}
The Gaussian filter smooths the data while preserving general brightness structure; the center pixel becomes a high-confidence weighted local average.

\subsection{Example Problem 2: Median Filter for Impulse Noise}

\subsubsection*{Problem}
Apply the $3 \times 3$ median filter to the following neighborhood:
\[
\{200, 205, 0,\; 198, 202, 210,\; 0, 199, 203\}
\]
Assume the zeros are salt-and-pepper noise. Compute the filtered output.

\subsubsection*{Solution}
Sort the values:
\[
0, 0, 198, 199, 200, 202, 203, 205, 210
\]
Median = 200.

\subsubsection*{Interpretation}
The median filter removes impulse noise while preserving edges, unlike averaging filters which would smear pixel intensities.
\section{Common Filters Used in Image Processing}

This section summarizes the most important spatial-domain filters used in digital image processing.  
For each filter, we provide:  
\begin{itemize}
\item The purpose and intuition  
\item Mathematical/matrix form  
\item Image-processing–specific usage  
\end{itemize}

All filters operate on local neighborhoods (usually $3 \times 3$ or $5 \times 5$).

\subsection{Smoothing (Low-Pass) Filters}

\subsubsection{Averaging Filter (Mean Filter)}
Reduces noise by computing the mean of pixel intensities.

\paragraph{3 × 3 Mask:}
\[
H = \frac{1}{9}
\begin{bmatrix}
1 & 1 & 1\\
1 & 1 & 1\\
1 & 1 & 1
\end{bmatrix}
\]

\paragraph{Usage:}
Noise reduction, rough smoothing. Blurs edges significantly.

\subsubsection{Weighted Averaging Filter}
Gives more weight to central pixels.

\[
H = \frac{1}{16}
\begin{bmatrix}
1 & 2 & 1\\
2 & 4 & 2\\
1 & 2 & 1
\end{bmatrix}
\]

\paragraph{Usage:}
Better smoothing with less edge distortion than the simple mean filter.

\subsubsection{Gaussian Filter}
Approximates a Gaussian blur, optimal for noise reduction without ringing.

\[
H = \frac{1}{16}
\begin{bmatrix}
1 & 2 & 1\\
2 & 4 & 2\\
1 & 2 & 1
\end{bmatrix}
\]

(This is a discrete approximation for $\sigma=1$.)

\paragraph{Usage:}
Reduces Gaussian noise; preserves edge contours better than averaging.

\subsection{Sharpening (High-Pass) Filters}

Sharpening filters enhance edges by emphasizing high-frequency components.

\subsubsection{Laplacian Filter}
Isotropic second-derivative operator.

\paragraph{4-neighbor version:}
\[
H =
\begin{bmatrix}
0 & 1 & 0\\
1 & -4 & 1\\
0 & 1 & 0
\end{bmatrix}
\]

\paragraph{8-neighbor version:}
\[
H =
\begin{bmatrix}
1 & 1 & 1\\
1 & -8 & 1\\
1 & 1 & 1
\end{bmatrix}
\]

\paragraph{Usage:}
Edge enhancement, image sharpening, feature extraction.

\subsubsection{Sharpening via Unsharp Masking}
\[
g = f - f_{LP}
\]

\subsubsection{High-Boost Filtering}
\[
g = A f - f_{LP}, \quad A > 1
\]

Matrix-based implementation uses any smoothing filter for $f_{LP}$.

\subsection{Gradient-Based Filters}

Approximations to the first derivative for edge detection.

\subsubsection{Sobel Filters}

\paragraph{Horizontal Sobel ($G_x$):}
\[
G_x =
\begin{bmatrix}
-1 & 0 & 1\\
-2 & 0 & 2\\
-1 & 0 & 1
\end{bmatrix}
\]

\paragraph{Vertical Sobel ($G_y$):}
\[
G_y =
\begin{bmatrix}
-1 & -2 & -1\\
0 & \;\;0 & 0\\
1 & 2 & 1
\end{bmatrix}
\]

\paragraph{Usage:}
Edge detection, gradient magnitude computation.

\subsubsection{Prewitt Filters}

\paragraph{Horizontal:}
\[
P_x =
\begin{bmatrix}
-1 & 0 & 1\\
-1 & 0 & 1\\
-1 & 0 & 1
\end{bmatrix}
\]

\paragraph{Vertical:}
\[
P_y =
\begin{bmatrix}
-1 & -1 & -1\\
0 & \;\;0 & 0\\
1 & 1 & 1
\end{bmatrix}
\]

\paragraph{Usage:}
Simple edge detection; less smoothing than Sobel.

\subsection{Nonlinear Filters}

\subsubsection{Median Filter}
Replaces each pixel with the median of the neighborhood.

\[
g[m,n] = \text{median} \{ f[i,j] \, | \, (i,j) \in \mathcal{N}_{m,n} \}
\]

\paragraph{Usage:}
Excellent for salt-and-pepper noise removal; preserves edges.

\subsubsection{Max Filter}
\[
g[m,n] = \max \{ f[i,j] \, | \, (i,j) \in \mathcal{N}_{m,n} \}
\]

\paragraph{Usage:}
Extracts bright regions; used in morphological dilation as well.

\subsubsection{Min Filter}
\[
g[m,n] = \min \{ f[i,j] \, | \, (i,j) \in \mathcal{N}_{m,n} \}
\]

\paragraph{Usage:}
Extracts dark regions; related to morphological erosion.

\subsubsection{Midpoint Filter}
\[
g[m,n] = \frac{1}{2}(\max + \min)
\]

\paragraph{Usage:}
Useful for uniform noise reduction.

\subsection{Bandpass and Bandreject Filters (Spatial Approximations)}

\subsubsection{High-Pass Approximation}
\[
H_{HP} =
\begin{bmatrix}
-1 & -1 & -1\\
-1 & 8 & -1\\
-1 & -1 & -1
\end{bmatrix}
\]

\paragraph{Usage:}
Emphasizes edges and fine detail.

\subsubsection{Low-Pass Approximation}
Same as averaging or Gaussian.

\subsubsection{Bandpass}
\[
H_{BP} = H_{HP} * H_{LP}
\]

\paragraph{Usage:}
Texture extraction, noise isolation.

\subsection{Directional Filters}

\subsubsection{Horizontal Line Detector}
\[
H =
\begin{bmatrix}
-1 & -1 & -1\\
2 & 2 & 2\\
-1 & -1 & -1
\end{bmatrix}
\]

\subsubsection{Vertical Line Detector}
\[
H =
\begin{bmatrix}
-1 & 2 & -1\\
-1 & 2 & -1\\
-1 & 2 & -1
\end{bmatrix}
\]

\paragraph{Usage:}
Extraction of oriented features (roads, edges, lines).

\subsection{Summary}

\begin{itemize}
\item Smoothing filters reduce noise but blur edges.
\item Sharpening filters enhance edges and fine detail.
\item Gradient filters detect edges using first derivatives.
\item Nonlinear filters (median, min/max) are essential for salt-and-pepper noise.
\item Directional filters detect oriented structures.
\end{itemize}

\subsection{Summary}

\begin{itemize}
\item Filtering modifies an image using either linear or nonlinear operations.
\item Smoothing filters reduce noise; sharpening filters enhance edges.
\item Median filters are robust against impulse noise.
\item Filters can be implemented in spatial or frequency domains depending on the application.
\end{itemize}

\section{Fourier Transform}

\subsection{Theory}

The Fourier Transform is one of the most powerful tools in image processing.  
It decomposes an image into its sinusoidal frequency components, representing it in the frequency domain.  

High frequencies correspond to:
\begin{itemize}
\item edges,
\item fine details,
\item rapid intensity changes.
\end{itemize}

Low frequencies correspond to:
\begin{itemize}
\item smooth variations,
\item illumination gradients,
\item background components.
\end{itemize}

Fourier methods are essential for:
\begin{itemize}
\item filtering (low-pass, high-pass, bandpass),
\item image restoration,
\item compression,
\item spectral analysis,
\item PSF estimation.
\end{itemize}

\subsection{Mathematical Formulation}

\subsubsection{Continuous 2D Fourier Transform}

For a continuous image $f(x,y)$:
\begin{equation}
F(u,v) = \int_{-\infty}^{\infty}\int_{-\infty}^{\infty}
f(x,y)\, e^{-j2\pi(ux + vy)} \, dx \, dy
\end{equation}

Inverse transform:
\begin{equation}
f(x,y) = \int_{-\infty}^{\infty}\int_{-\infty}^{\infty}
F(u,v)\, e^{j2\pi(ux + vy)} \, du \, dv
\end{equation}

\subsubsection{Discrete 2D Fourier Transform (DFT)}

Digital images use the discrete form:
\begin{equation}
F(u,v) =
\sum_{x=0}^{M-1} \sum_{y=0}^{N-1}
f(x,y)e^{-j2\pi \left( \frac{ux}{M} + \frac{vy}{N} \right)}
\end{equation}

Inverse DFT:
\begin{equation}
f(x,y) = \frac{1}{MN}
\sum_{u=0}^{M-1} \sum_{v=0}^{N-1}
F(u,v)e^{j2\pi \left( \frac{ux}{M} + \frac{vy}{N} \right)}
\end{equation}

\paragraph{Magnitude and Phase:}
\[
|F(u,v)| = \sqrt{\Re^2 + \Im^2}, \qquad
\phi(u,v) = \tan^{-1}\left(\frac{\Im}{\Re}\right)
\]

\textbf{Key fact:} Image structure comes primarily from the \textbf{phase}.

\subsection{Key Properties of the Fourier Transform}

\subsubsection{Linearity}
\[
\mathcal{F}\{af + bg\} = aF + bG
\]

\subsubsection{Shift Property}
A spatial shift introduces a phase shift:
\[
f(x-x_0, y-y_0)
\; \xleftrightarrow{\mathcal{F}} \;
e^{-j2\pi(ux_0/M + vy_0/N)}F(u,v)
\]

\subsubsection{Convolution Theorem}
\[
f * h \;\xleftrightarrow{\mathcal{F}}\; FH
\]

\paragraph{Most important result in image filtering.}

\subsubsection{Correlation Theorem}
\[
f \circ h \;\xleftrightarrow{\mathcal{F}}\; F H^*
\]

\subsubsection{Separability}
If
\[
H(u,v) = H_1(u)H_2(v),
\]
then filtering can be performed using 1D transforms (computationally efficient).

\subsubsection{Symmetry for Real Images}
For real-valued images:
\[
F(-u,-v) = F^*(u,v)
\]

\subsubsection{Centering the Fourier Transform}
To move the DC component to the center:
\[
f_c(x,y) = f(x,y)(-1)^{x+y}
\]

\subsection{Interpretation in Image Processing}

\begin{itemize}
\item Low frequencies = illumination + smooth background.
\item High frequencies = edges, fine textures.
\item Filtering in frequency domain is multiplication:
\[
G(u,v) = H(u,v)F(u,v)
\]
\item Noise removal, deblurring, and PSF inversion rely heavily on Fourier analysis.
\item Compression (JPEG) uses frequency decomposition (via DCT).
\end{itemize}

\subsection{Example Problem 1: Fourier Transform of a Separable Image}

\subsubsection*{Problem}
Let
\[
f(x,y) = a(x)b(y)
\]
where $a(x) = [1, 2]$ and $b(y) = [3, 1]$ (1D signals).  
Compute the 2D DFT $F(u,v)$.

\subsubsection*{Solution}

Since $f(x,y)$ is separable:
\[
F(u,v) = A(u) B(v)
\]

Compute 1D DFT of $a(x)$:
\[
A(0) = 1 + 2 = 3
\]
\[
A(1) = 1 - 2 = -1
\]

Compute 1D DFT of $b(y)$:
\[
B(0) = 3 + 1 = 4
\]
\[
B(1) = 3 - 1 = 2
\]

Then the outer product:
\[
F =
\begin{bmatrix}
A(0)B(0) & A(0)B(1)\\
A(1)B(0) & A(1)B(1)
\end{bmatrix}
=
\begin{bmatrix}
3\cdot 4 & 3\cdot 2\\
-1\cdot 4 & -1\cdot 2
\end{bmatrix}
\]

Thus:
\[
F =
\begin{bmatrix}
12 & 6\\
-4 & -2
\end{bmatrix}
\]

\subsubsection*{Interpretation}
The entire 2D spectrum can be computed from two 1D transforms — illustrating separability.

\subsection{Example Problem 2: Effect of Image Shift on Fourier Transform}

\subsubsection*{Problem}
If
\[
g(x,y) = f(x-2, y-3),
\]
derive the Fourier transform $G(u,v)$ in terms of $F(u,v)$.

\subsubsection*{Solution}

Using the shift property:
\[
g(x,y) \xleftrightarrow{\mathcal{F}}
e^{-j2\pi\left(\frac{2u}{M} + \frac{3v}{N}\right)} F(u,v)
\]

Thus:
\[
G(u,v) = F(u,v) e^{-j2\pi\left(\frac{2u}{M} + \frac{3v}{N}\right)}
\]

\subsubsection*{Interpretation}
A spatial shift does **not** change magnitude, only phase — demonstrating why phase contains structural information.

\subsection{Summary}

\begin{itemize}
\item The Fourier Transform converts images into frequency components.
\item High and low frequencies correspond to edges and smooth regions, respectively.
\item The convolution theorem enables efficient frequency-domain filtering.
\item Shifts in space produce phase changes in the frequency domain.
\item Fourier analysis is foundational for restoration, enhancement, compression, and PSF modeling.
\end{itemize}

\section{Inverse Fourier Transform}

\subsection{Theory}

The Inverse Fourier Transform (IFT) reconstructs an image from its frequency domain representation.  
While the Fourier Transform decomposes an image into sinusoidal components, the inverse transform recombines them into the original spatial-domain image.

In image processing, the IFT is essential for:
\begin{itemize}
\item reconstructing filtered images (after applying $G(u,v) = H(u,v)F(u,v)$),
\item image restoration and deblurring,
\item reconstruction in MRI, CT, and other spectral imaging applications,
\item transitioning between spatial and frequency domain algorithms,
\item verifying linear shift-invariant (LSI) system models.
\end{itemize}

The forward and inverse transforms form a dual pair, allowing lossless switching between domains (for ideal sampling conditions).

\subsection{Mathematical Formulation}

\subsubsection{Continuous 2D Inverse Fourier Transform}

Given the Fourier Transform:
\[
F(u,v) = \iint f(x,y)e^{-j2\pi(ux+vy)} \, dx\,dy
\]

The inverse is:
\[
f(x,y) = \iint F(u,v)e^{j2\pi(ux+vy)} \, du\,dv
\]

The positive sign in the exponential restores the spatial-domain signal.

\subsubsection{Discrete 2D Inverse Fourier Transform (IDFT)}

For digital images, the inverse is:
\begin{equation}
f(x,y) = \frac{1}{MN}
\sum_{u=0}^{M-1} \sum_{v=0}^{N-1}
F(u,v)e^{j2\pi \left( \frac{ux}{M} + \frac{vy}{N} \right)}
\end{equation}

\paragraph{Key Points:}
\begin{itemize}
\item The normalization factor $\frac{1}{MN}$ ensures correct scaling.
\item When $F(u,v)$ is complex, the reconstructed $f(x,y)$ is real if the input is real.
\item Image-processing software (MATLAB, NumPy, OpenCV) uses the same definition.
\end{itemize}

\subsection{Properties of the Inverse Fourier Transform}

\subsubsection{1. Duality}
If
\[
\mathcal{F}\{f(x,y)\} = F(u,v),
\]
then
\[
\mathcal{F}^{-1}\{F(x,y)\} = f(-u,-v)
\]

\subsubsection{2. Uniqueness}
If $f(x,y)$ is real and finite, then $F(u,v)$ uniquely determines the image.

\subsubsection{3. Hermitian Symmetry}
For real images $f(x,y)$:
\[
F(-u,-v) = F^*(u,v)
\]
This ensures IFT yields a real-valued image.

\subsubsection{4. Energy Conservation (Parseval’s Theorem)}
\[
\sum_{x,y} |f(x,y)|^2 = \frac{1}{MN} \sum_{u,v} |F(u,v)|^2
\]
Useful for comparing filtered vs. noisy images.

\subsubsection{5. Circular Convolution}
DFT and IDFT obey circular convolution:
\[
f * h \quad \leftrightarrow \quad FH
\]
Padding is needed to avoid wraparound artifacts.

\subsection{Interpretation in Image Processing}

\begin{itemize}
\item Every frequency-domain filter must be returned to the spatial domain using the IFT.
\item Highpass, lowpass, bandpass, and sharpening filters are implemented as:
\[
g(x,y) = \mathcal{F}^{-1}\{H(u,v)F(u,v)\}
\]
\item Image restoration (inverse filtering, Wiener filtering) relies entirely on this process.
\item After applying FFT-based convolution, the IFT retrieves the processed image:
\[
g = \text{IFFT}(\text{FFT}(f)\cdot \text{FFT}(h))
\]
\end{itemize}

\subsection{Example Problem 1: Recovering Spatial Image from Simple Spectrum}

\subsubsection*{Problem}

Given the 2D DFT:
\[
F =
\begin{bmatrix}
12 & 0\\
0 & -4
\end{bmatrix}
\]
for a $2\times2$ image, compute the inverse DFT to obtain $f(x,y)$.

\subsubsection*{Solution}

Using the IDFT:
\[
f(x,y) = \frac{1}{4}
\sum_{u=0}^{1}\sum_{v=0}^{1}
F(u,v)e^{j2\pi\left(\frac{ux}{2} + \frac{vy}{2}\right)}
\]

Compute all four values:

\[
f(0,0) = \frac{1}{4}(12 + 0 + 0 - 4) = 2
\]

\[
f(0,1) = \frac{1}{4}(12 - 0 + 0 + 4) = 4
\]

\[
f(1,0) = \frac{1}{4}(12 + 0 - 0 + 4) = 4
\]

\[
f(1,1) = \frac{1}{4}(12 - 0 - 0 - 4) = 2
\]

Thus:
\[
f =
\begin{bmatrix}
2 & 4\\
4 & 2
\end{bmatrix}
\]

\subsubsection*{Interpretation}
Low frequencies dominate → smooth spatial pattern with symmetric values.

\subsection{Example Problem 2: IFT After Filtering}

\subsubsection*{Problem}
A $2\times2$ image has DFT:
\[
F =
\begin{bmatrix}
20 & -4\\
-4 & 0
\end{bmatrix}
\]
A lowpass filter keeps only $F(0,0)$ and sets others to zero.  
Compute the spatial-domain output.

\subsubsection*{Solution}

Filtered spectrum:
\[
G =
\begin{bmatrix}
20 & 0\\
0 & 0
\end{bmatrix}
\]

Apply IDFT:
\[
g(x,y) = \frac{1}{4} ( 20 )
\]

Thus all pixels become:
\[
g = 5
\begin{bmatrix}
1 & 1\\
1 & 1
\end{bmatrix}
\]

\subsubsection*{Interpretation}
Only the DC component remains → completely smoothed (constant) image.

\subsection{Summary}

\begin{itemize}
\item The inverse Fourier Transform reconstructs an image from its spectral representation.
\item Filtering in the frequency domain always requires the IFT to obtain spatial results.
\item Properties such as Hermitian symmetry ensure real-valued outputs.
\item The IFT is essential for restoration, filtering, and frequency-based compression.
\end{itemize}

\section{Quantization}

\subsection{Theory}

Quantization is the process of mapping a continuous or large-range set of pixel intensity values into a finite set of discrete levels.  
In digital image processing, an image must be quantized in intensity to convert it from an analog signal into a digital representation.

An image with bit depth $b$ has:
\[
L = 2^b \quad \text{quantization levels}
\]

Quantization introduces:
\begin{itemize}
\item \textbf{rounding error (quantization error)}
\item \textbf{loss of fine detail}
\item \textbf{contouring artifacts} (visible steps in smooth gradients)
\end{itemize}

\paragraph{Sources of quantization:}
\begin{itemize}
\item Analog-to-digital conversion (ADC)
\item Digital compression algorithms
\item Display systems (e.g., 8-bit grayscale monitors)
\end{itemize}

\subsection{Mathematical Formulation}

Let the continuous intensity range be:
\[
[I_{\min}, I_{\max}]
\]

For uniform quantization with $L$ levels, the quantization step is:
\[
\Delta = \frac{I_{\max} - I_{\min}}{L - 1}
\]

A pixel value $r$ is mapped to:
\[
Q(r) = \Delta \cdot \text{round}\left(\frac{r-I_{\min}}{\Delta}\right) + I_{\min}
\]

\subsubsection{Quantization Error}

The quantization error is:
\[
e_q = r - Q(r)
\]
For uniform quantizers, it is bounded by:
\[
-\frac{\Delta}{2} \leq e_q \leq \frac{\Delta}{2}
\]

Assuming a uniform distribution of error:
\[
\sigma_{e}^{2} = \frac{\Delta^2}{12}
\]

\subsubsection{Signal-to-Noise Ratio (SNR)}

For an image with variance $\sigma_f^2$, the quantization SNR is:
\[
\text{SNR}_q = 10\log_{10}\left(\frac{\sigma_f^2}{\sigma_e^2}\right)
\]

\textbf{As $b$ increases, $\Delta$ decreases → better image quality.}

\subsection{Types of Quantizers}

\subsubsection{Uniform Quantization}
\[
Q(r) = \Delta \cdot \text{round}\left(\frac{r}{\Delta}\right)
\]

Used in:
\begin{itemize}
\item Standard grayscale imaging (8-bit, 12-bit, 16-bit)
\item JPEG and DCT-based quantizers (uniform in transformed domain)
\end{itemize}

\subsubsection{Non-Uniform Quantization}
More levels are allocated to frequently occurring intensity values.

Optimal quantizers minimize mean square error:
\[
Q^* = \arg\min_{Q} \; E[(f - Q(f))^2]
\]

Examples:
\begin{itemize}
\item Lloyd-Max quantizer
\item Mu-law / A-law quantization (for compressed sensing or HDR imaging)
\end{itemize}

\subsection{Effects of Quantization in Imaging}

\subsubsection{False Contouring}
Occurs when $L$ is too small.  
Gradients appear as visible step edges (e.g., sky in 4-bit images).

\subsubsection{Loss of Detail}
Small variations within a region are merged into the same quantization bin.

\subsubsection{Quantization Noise in Smooth Areas}
Low-variance regions show more visible quantization noise.

\paragraph{Mitigation Strategies:}
\begin{itemize}
\item Increase bit depth
\item Apply dithering
\item Use nonlinear quantization
\end{itemize}

\subsection{Example Problem 1: Quantization Error and SNR}

\subsubsection*{Problem}

An image has pixel intensities in $[0, 255]$.  
It is quantized to 4 bits.  
Compute:

\begin{enumerate}
\item Quantization step $\Delta$
\item Maximum quantization error
\item Quantization noise variance
\item SNR if the image has variance $\sigma_f^2 = 1200$
\end{enumerate}

\subsubsection*{Solution}

Number of levels:
\[
L = 2^4 = 16
\]

(1) Quantization step:
\[
\Delta = \frac{255 - 0}{16 - 1} = \frac{255}{15} = 17
\]

(2) Max error:
\[
|e_q| \leq \frac{\Delta}{2} = 8.5
\]

(3) Noise variance:
\[
\sigma_e^2 = \frac{\Delta^2}{12} = \frac{17^2}{12}
= \frac{289}{12} = 24.08
\]

(4) SNR:
\[
\text{SNR} = 10 \log_{10} \left(\frac{1200}{24.08}\right)
= 10\log_{10}(49.8)
= 10 \cdot 1.697 = 16.97 \text{ dB}
\]

\subsubsection*{Interpretation}

A 4-bit image shows significant quantization noise—SNR of only 17 dB is poor for imaging.

\subsection{Example Problem 2: False Contouring Analysis}

\subsubsection*{Problem}

A smooth gradient image is quantized from 256 gray levels to 8 levels.  
Determine:

\begin{enumerate}
\item The quantization step size.
\item Whether false contouring will be visually noticeable.
\end{enumerate}

\subsubsection*{Solution}

\[
L = 8
\]

\[
\Delta = \frac{255}{7} \approx 36.4
\]

This means any gradual change smaller than 36 gray levels maps to the same output.

Human visual system detects changes as low as 2–3 gray levels, so 36-level jumps are extremely noticeable.

\subsubsection*{Conclusion}

Strong false contouring will be visible in sky regions, medical ultrasound scans, or any smoothly varying image.

\subsection{Summary}

\begin{itemize}
\item Quantization reduces continuous intensity values to discrete levels.
\item It introduces quantization error, modeled as uniform noise.
\item Low bit-depth images show false contouring and noise.
\item SNR improves by roughly 6 dB for every additional quantization bit.
\end{itemize}

\section{Histogram Equalization}

\subsection{Theory}

Histogram Equalization (HE) is a spatial domain technique used to improve global image contrast by redistributing the pixel intensities so that the output image has a \emph{uniform} histogram.  
Unlike simple contrast stretching, HE is automatic and data-driven.

The method comes directly from probability theory:  
If we apply a transformation based on the cumulative distribution function (CDF),  
\[
s = T(r)
\]
then the transformed random variable \(s\) becomes \textbf{uniformly distributed} in \([0,1]\).

Histogram equalization applies this idea to discrete image intensities.

\subsection{Mathematical Derivation (Based on Gonzalez and Woods)}

Let an image have gray levels:
\[
r_k,\quad k = 0,1,\ldots,L-1
\]
and a histogram giving pixel counts \(n_k\).  

Define the probability of gray level:
\[
p_r(r_k) = \frac{n_k}{MN}
\]
where \(M \times N\) is the image size.

\subsubsection{1. Transformation Function}

The histogram equalization transform is:
\[
s_k = T(r_k) = (L-1) \sum_{j=0}^{k} p_r(r_j)
\]
This is simply the CDF scaled to the range \([0, L-1]\).

\begin{itemize}
\item If many pixels are in a low-intensity region, the CDF rises quickly → stretches the dark region.
\item If few pixels appear in bright regions, the CDF rises slowly → compresses bright regions.
\end{itemize}

Thus contrast is enhanced by redistributing pixel intensities.

\subsubsection{2. Why Histogram Equalization Works}

For a continuous function, if:
\[
s = T(r)
\]
where \(T\) is strictly increasing and differentiable, and:
\[
T(r) = \int_{0}^{r} p_r(\omega) \, d\omega,
\]
then it can be shown that:
\[
p_s(s) = 1
\]
for \(s \in [0,1]\).  

Meaning: **the output distribution becomes uniform.**

In discrete digital images, uniformity is approximate—not exact—but the contrast enhancement effect remains very strong.

\subsection{Algorithmic Steps (Discrete Implementation)}

Given an image with intensity range \([0, L-1]\):

\begin{enumerate}
\item Compute histogram \(h(r_k)\)
\[
h(r_k) = n_k
\]
\item Normalize to obtain PMF
\[
p_r(r_k) = \frac{n_k}{MN}
\]
\item Compute CDF
\[
C(k) = \sum_{j=0}^{k} p_r(r_j)
\]
\item Multiply by \((L-1)\)
\[
s_k = (L-1) C(k)
\]
\item Apply mapping:  
Every pixel of value \(r_k\) is replaced by \( \lfloor s_k \rfloor \).
\end{enumerate}

\subsection{Interpretation}

Histogram equalization:
\begin{itemize}
\item Spreads out frequently occurring intensities
\item Increases global contrast
\item Enhances details in dark and bright regions simultaneously
\item Is fully automatic — no parameters required
\end{itemize}

However:
\begin{itemize}
\item It may over-enhance noise in homogeneous regions.
\item It may produce unnatural brightness for images intended for human interpretation (e.g., medical).
\end{itemize}

\subsection{Local Histogram Equalization (Adaptive HE and CLAHE)}

Global HE uses the histogram of the entire image → fails when illumination is nonuniform.

\subsubsection{Local HE}
Break image into regions (e.g., \(8 \times 8\) blocks):
\begin{enumerate}
\item Perform HE on each block.
\item Interpolate block boundaries to avoid artifacts.
\end{enumerate}

\subsubsection{CLAHE (Contrast Limited AHE)}
To prevent noise amplification:
\begin{itemize}
\item Clip histogram at a threshold (the "clip limit").
\item Redistribute clipped counts uniformly.
\item Compute CDF and equalize.
\end{itemize}

CLAHE is widely used in:
\begin{itemize}
\item Medical imaging (X-ray, CT)
\item Night-vision enhancement
\item Low-light photography
\end{itemize}

\subsection{Example Problem 1: Manual Histogram Equalization}

\subsubsection*{Problem}

Consider a \(3 \times 3\) image with intensity levels \(0,1,2,3\):  
\[
f =
\begin{bmatrix}
0 & 1 & 1 \\
2 & 3 & 3 \\
3 & 2 & 1
\end{bmatrix}
\]

Perform histogram equalization.

\subsubsection*{Solution}

\paragraph{1. Histogram}

\[
h(0)=1,\; h(1)=3,\; h(2)=2,\; h(3)=3
\]

\[
MN = 9,\quad L=4
\]

\paragraph{2. PMF}

\[
p_r = \{1/9,\; 3/9,\; 2/9,\; 3/9\}
\]

\paragraph{3. CDF}

\[
C(0)=1/9,\qquad
C(1)=4/9,\qquad
C(2)=6/9,\qquad
C(3)=9/9
\]

\paragraph{4. Mapping}

\[
s_k = (L-1) C(k) = 3C(k)
\]

\[
s_0 = 0.33 \to 0
\]
\[
s_1 = 1.33 \to 1
\]
\[
s_2 = 2.00 \to 2
\]
\[
s_3 = 3.00 \to 3
\]

Thus the mapping is:
\[
0 \to 0,\quad 1 \to 1,\quad 2 \to 2,\quad 3 \to 3
\]

Since the original histogram already had reasonably spread intensities, equalization does not change output much.

\subsubsection*{Result}

\[
g = f
\]

\subsection{Example Problem 2: Histogram Equalization on a Low-Contrast Image}

\subsubsection*{Problem}

A small low-contrast image has intensities in \(\{50, 52, 55\}\).  
Its histogram is:

\[
h(50)=5,\; h(52)=2,\; h(55)=1
\]

Perform histogram equalization assuming 8-bit output \((L = 256)\).

\subsubsection*{Solution}

Compute PMF:

\[
p_r(50)=5/8, \quad p_r(52)=2/8, \quad p_r(55)=1/8
\]

CDF:

\[
C(50)=5/8,\quad C(52)=7/8,\quad C(55)=1
\]

Mapping to 0–255 range:

\[
s(50) = 255(5/8) = 159
\]
\[
s(52) = 255(7/8) = 223
\]
\[
s(55) = 255(1) = 255
\]

\paragraph{Interpretation}

The intensities \(\{50,52,55\}\) are now remapped to \(\{159,223,255\}\):

\begin{itemize}
\item Very large contrast expansion  
\item Highlights previously invisible details  
\item Corrects for the low dynamic range of the original image
\end{itemize}

\subsection{Summary}

\begin{itemize}
\item Histogram Equalization uses the CDF as a transformation function.
\item It expands frequently occurring intensity ranges → increases contrast.
\item It is automatic but global; local variants (AHE, CLAHE) improve robustness.
\item HE is fundamental for preprocessing in vision, medical imaging, and remote sensing.
\end{itemize}
\section{Degradation and the Point Spread Function (PSF)}

\subsection{Introduction}

Image degradation refers to any process that causes the captured image to differ from the ideal scene.  
Examples include:
\begin{itemize}
\item motion blur,
\item atmospheric turbulence,
\item out-of-focus blur,
\item sensor imperfections and noise.
\end{itemize}

Mathematically, degradation is modeled as a linear, shift-invariant (LSI) system characterized by a \textbf{Point Spread Function (PSF)}.

\subsection{Degradation Model}

The most general degradation model in the spatial domain is:
\begin{equation}
g(x,y) = h(x,y) * f(x,y) + \eta(x,y)
\end{equation}

Where:
\begin{itemize}
\item $f(x,y)$ : original image  
\item $h(x,y)$ : point spread function (PSF)  
\item $g(x,y)$ : observed degraded image  
\item $\eta(x,y)$ : additive noise  
\item $*$ : 2D convolution  
\end{itemize}

Taking the Fourier transform yields:
\begin{equation}
G(u,v) = H(u,v)F(u,v) + N(u,v)
\end{equation}

Where $H(u,v)$ is the \textbf{Optical Transfer Function (OTF)}, i.e., the Fourier transform of the PSF.

\subsection{Point Spread Function (PSF)}

The PSF describes the response of an imaging system to a point source of light.  
In ideal imaging, the PSF is a delta function:
\[
h(x,y) = \delta(x,y)
\]
meaning no blurring.

Real-world systems spread light, producing a blur kernel.

\subsubsection{Properties of PSF}

\begin{itemize}
\item Spatial domain representation of blur
\item Must be non-negative
\item Must integrate to 1 (energy preservation)
\[
\iint h(x,y) \, dx\,dy = 1
\]
\item Determines the frequency response:
\[
H(u,v) = \mathcal{F}\{h(x,y)\}
\]
\end{itemize}

\subsection{Common PSF Models in Imaging}

\subsubsection{1. Motion Blur PSF (Uniform Linear Motion)}

If camera or object moves uniformly during exposure, the PSF is:
\[
h(x,y) = \frac{1}{L} \text{sinc-like blur along motion direction}
\]

More precisely, for motion of length \( L \) at angle \( \theta \):
\[
h(x,y) = 
\begin{cases}
\frac{1}{L}, & 0 \le x\cos\theta + y \sin\theta \le L \\
0, & \text{otherwise}
\end{cases}
\]

Its frequency response:
\[
H(u,v) = \frac{\sin(\pi (u\cos\theta + v\sin\theta)L)}
{\pi (u\cos\theta + v\sin\theta)L}
\]

\subsubsection{2. Out-of-Focus Blur (Defocus Blur)}

A circular aperture gives:
\[
h(x,y) =
\begin{cases}
\frac{1}{\pi R^2}, & x^2 + y^2 \le R^2 \\
0, & \text{otherwise}
\end{cases}
\]

OTF:
\[
H(u,v) = \frac{2 J_1 (2\pi R D)}{2\pi R D}
\]
where \( J_1 \) is the Bessel function and \( D=\sqrt{u^2+v^2} \).

\subsubsection{3. Atmospheric Turbulence Blur}

Modeled using:
\[
H(u,v) = e^{-k(u^2+v^2)^{5/6}}
\]

Used in satellite and astronomy imaging.

\subsubsection{4. Gaussian Blur PSF}

For optical scattering:
\[
h(x,y) = \frac{1}{2\pi\sigma^2} \exp\left( -\frac{x^2 + y^2}{2\sigma^2} \right)
\]

Fourier transform:
\[
H(u,v) = \exp\left( -2\pi^2\sigma^2(u^2+v^2) \right)
\]

This is widely used due to its simplicity and physical relevance.

\subsection{Noise in the Degradation Model}

Noise is typically independent of blur.  
Common noise distributions:
\begin{itemize}
\item Gaussian noise (sensor/thermal noise)
\item Salt-and-pepper noise (bit errors)
\item Poisson noise (photon counting)
\end{itemize}

Noise makes restoration far more challenging.

\subsection{Why PSF Matters in Restoration}

All restoration methods require knowledge or estimation of the PSF.

\begin{itemize}
\item \textbf{Inverse filtering:}
\[
\hat{F}(u,v) = \frac{G(u,v)}{H(u,v)}
\]
Requires exact PSF.

\item \textbf{Wiener filtering:}
\[
\hat{F}(u,v) = \frac{H^*(u,v)}
{|H(u,v)|^2 + S_\eta(u,v)/S_f(u,v)} G(u,v)
\]

\item \textbf{Constrained least squares (CLS).}

\item \textbf{Blind deconvolution:}  
Estimate \( h \) and \( f \) simultaneously.
\end{itemize}

\subsection{Example Problem 1: Motion Blur PSF and OTF}

\subsubsection*{Problem}

A camera moves horizontally by 10 pixels during exposure.  
Assume motion blur PSF:
\[
h(x,y) = 
\begin{cases}
\frac{1}{10}, & 0 \le x \le 10, \; y=0\\
0, & \text{otherwise}
\end{cases}
\]

Find the frequency response \( H(u,v) \).

\subsubsection*{Solution}

By definition:
\[
H(u,v) = \int_{0}^{10} \frac{1}{10} e^{-j2\pi ux} dx
\]

Compute integral:
\[
H(u,v) = \frac{1}{10} \left[ \frac{e^{-j2\pi ux}}{-j2\pi u} \right]_{0}^{10}
\]

\[
H(u,v) = \frac{1}{10} \cdot \frac{1}{j2\pi u}
\left(1 - e^{-j20\pi u}\right)
\]

Using Euler form:
\[
H(u,v) = \frac{\sin(10\pi u)}{10\pi u}
\]

\[
\boxed{
H(u,v) = \text{sinc}(10u)
}
\]

This is the classical motion-blur frequency response.

\subsection{Example Problem 2: Gaussian PSF Degradation}

\subsubsection*{Problem}

An imaging system has Gaussian PSF:
\[
h(x,y)=\frac{1}{2\pi\sigma^2}e^{-(x^2+y^2)/(2\sigma^2)}
\]
Find the OTF and determine at what frequency the magnitude drops to 0.5.

Let \( \sigma = 2 \).

\subsubsection*{Solution}

The OTF of a Gaussian is:
\[
H(u,v) = e^{-2\pi^2\sigma^2 (u^2+v^2)}
\]

Set:
\[
|H| = 0.5
\]

\[
0.5 = e^{-2\pi^2(4)(u^2+v^2)}
\]

Take logs:
\[
\ln 0.5 = -8\pi^2 (u^2+v^2)
\]

\[
u^2+v^2 = \frac{\ln 2}{8\pi^2}
\]

Numerically:
\[
u^2+v^2 \approx \frac{0.693}{78.96} = 0.00878
\]

Thus the radius in frequency domain at half-power is:
\[
D = \sqrt{0.00878} \approx 0.094
\]

Meaning the system strongly suppresses frequencies above 0.1 cycles/pixel.

\subsection{Summary}

\begin{itemize}
\item The PSF models blur in an imaging system.
\item Degradation model:
\[
g = h * f + \eta
\]
\item PSF’s Fourier transform is the OTF, essential in restoration.
\item Common PSFs model motion, defocus, atmospheric, and Gaussian blur.
\item Restoration methods (inverse, Wiener, CLS) depend on accurate PSF knowledge.
\end{itemize}

\section{Discrete Cosine Transform (DCT)}

\subsection{Introduction}

The Discrete Cosine Transform (DCT) is one of the most important transforms in image processing, especially in compression standards such as JPEG, MPEG, and H.26x.  
It provides excellent \textbf{energy compaction}, meaning most of the image energy is packed into a few low-frequency coefficients.

The DCT is closely related to the Fourier Transform, but uses only cosine basis functions, making it:

\begin{itemize}
\item real-valued (no imaginary components),
\item even-symmetric,
\item better suited for natural images,
\item less sensitive to boundary discontinuities.
\end{itemize}

The most commonly used version is the \textbf{2D DCT-II}, which is the mathematical core of JPEG.

\subsection{1D DCT Definition}

For a 1D signal $f(x)$ of length $N$:
\begin{equation}
F(u) = \alpha(u) \sum_{x=0}^{N-1} f(x)
\cos\left[\frac{\pi(2x+1)u}{2N}\right]
\end{equation}

Where:
\[
\alpha(u) =
\begin{cases}
\sqrt{\frac{1}{N}}, & u = 0\\[6pt]
\sqrt{\frac{2}{N}}, & u > 0
\end{cases}
\]

The inverse DCT (IDCT) is:
\begin{equation}
f(x) = \sum_{u=0}^{N-1} \alpha(u) F(u)
\cos\left[\frac{\pi(2x+1)u}{2N}\right]
\end{equation}

\subsection{2D DCT (Most Relevant for Images)}

For an $N \times N$ block (JPEG uses $N=8$):

\begin{equation}
F(u,v) = 
\alpha(u)\alpha(v)
\sum_{x=0}^{N-1}\sum_{y=0}^{N-1}
f(x,y)\cos\left[\frac{\pi(2x+1)u}{2N}\right]
\cos\left[\frac{\pi(2y+1)v}{2N}\right]
\end{equation}

The inverse transform is:

\begin{equation}
f(x,y) =
\sum_{u=0}^{N-1}\sum_{v=0}^{N-1}
\alpha(u)\alpha(v) F(u,v)
\cos\left[\frac{\pi(2x+1)u}{2N}\right]
\cos\left[\frac{\pi(2y+1)v}{2N}\right]
\end{equation}

\subsection{Interpretation of DCT Coefficients}

In the DCT domain:

\begin{itemize}
\item \(F(0,0)\) = DC coefficient (average intensity of the block).  
\item Low-frequency coefficients → smooth variations.  
\item High-frequency coefficients → edges, noise, details.
\end{itemize}

Natural images tend to concentrate most energy in low-frequency components.  
Thus high-frequency components can be quantized more aggressively → basis of JPEG compression.

\subsection{Why the DCT Compacts Energy Well}

Consider an $N \times N$ image block:

\begin{itemize}
\item Smooth regions → cosine basis aligns well → large values in DC and few low-frequency coefficients.
\item Sharp edges → energy spreads, but still more compact than Fourier.
\item Cosines reduce boundary discontinuity artifacts because they assume even symmetry.
\end{itemize}

\subsection{DCT in JPEG Compression}

JPEG applies DCT as follows:

\begin{enumerate}
\item Convert RGB to YCbCr.
\item Divide image into \(8\times8\) blocks.
\item Apply 2D DCT to each block.
\item Quantize each coefficient using quantization matrix \(Q(u,v)\):
\[
\hat{F}(u,v) = \text{round}\left( \frac{F(u,v)}{Q(u,v)} \right)
\]
\item Zig-zag scan (run-length encoding).
\item Entropy coding (Huffman or arithmetic coding).
\end{enumerate}

The DCT is thus the main step determining compression quality.

\subsection{Example Problem 1: Compute DCT of a Simple 1D Signal}

\subsubsection*{Problem}

Given a 1D signal of length 4:
\[
f(x) = [10,\; 10,\; 10,\; 10]
\]
Compute the DCT coefficients.

\subsubsection*{Solution}

Since the signal is constant:

\[
F(0) = \sqrt{\frac{1}{4}}(10+10+10+10) = \frac{1}{2}(40) = 20
\]

For \(u>0\), cosine terms integrate to zero due to orthogonality:

\[
F(1)= F(2)= F(3) = 0
\]

Thus:
\[
F = [20,\; 0,\; 0,\; 0]
\]

\subsubsection*{Interpretation}

All energy is in the DC component → perfect energy compaction.

\subsection{Example Problem 2: 2D DCT of an $2\times2$ Block}

\subsubsection*{Problem}

Compute the 2D DCT of:
\[
f =
\begin{bmatrix}
1 & 1\\
1 & 1
\end{bmatrix}
\]
Use $N=2$ and:
\[
\alpha(0)=\frac{1}{\sqrt{2}},\quad \alpha(1)=1
\]

\subsubsection*{Solution}

\paragraph{Compute $F(0,0)$ (DC term):}

\[
F(0,0) = \alpha(0)\alpha(0)
\sum f(x,y)
= \frac{1}{2}(4) = 2
\]

\paragraph{Compute $F(0,1)$:}

Cosine terms for $v=1$ alternate:

\[
F(0,1)=\alpha(0)\alpha(1)[1 - 1 + 1 - 1] = 0
\]

\paragraph{Compute $F(1,0)$:}

\[
F(1,0)=\alpha(1)\alpha(0)[1+1-1-1] = 0
\]

\paragraph{Compute $F(1,1)$:}

\[
F(1,1)=1\cdot 1 [1 - 1 -1 + 1] = 0
\]

Thus:
\[
F =
\begin{bmatrix}
2 & 0\\
0 & 0
\end{bmatrix}
\]

\subsubsection*{Interpretation}

Again, all energy is in the DC coefficient because the block is constant.

\subsection{Summary}

\begin{itemize}
\item The DCT uses cosine basis functions to transform an image block into frequency components.
\item Excellent for energy compaction → crucial for JPEG.
\item DC component represents average intensity; higher frequencies represent details.
\item High-frequency coefficients can be quantized more aggressively with minimal perceptual loss.
\end{itemize}

\section{Text-Based Data Compression}

\subsection{Introduction}

Data compression refers to representing information using fewer bits than the original representation.  
Compression is essential in storage, transmission, and multimedia systems (e.g., JPEG, PNG, MPEG).  

Even though this topic arises in the context of text compression, the principles apply directly to \emph{image compression}, where pixel intensities or transformed coefficients are treated as symbols.

\subsection{Types of Redundancy}

Compression works by removing redundancy. According to Gonzalez and Woods and classical information theory, there are three primary types:

\begin{enumerate}
\item \textbf{Coding redundancy}  
Some symbols use more bits than necessary.  
Example: fixed-length ASCII (8 bits) for 26 letters.

\item \textbf{Spatial (interpixel) redundancy}  
Pixels in images are correlated → predictable.  
Used by predictive coding, transform coding, LZ-family algorithms.

\item \textbf{Psychovisual redundancy}  
The human visual system (HVS) is less sensitive to high frequencies.  
Used in JPEG quantization.
\end{enumerate}

Lossless compression removes (1) and (2).  
Lossy compression additionally removes (3).

\subsection{Information Theory Basics}

\subsubsection{Entropy}

Let symbols have probabilities \(p_i\).  
The entropy of the source is:

\begin{equation}
H = - \sum_{i=1}^{N} p_i \log_2 p_i
\end{equation}

\paragraph{Properties:}
\begin{itemize}
\item \(H\) is the lower bound (in bits/symbol) of any lossless code.
\item More skewed distributions → lower entropy.
\end{itemize}

\subsubsection{Average Codeword Length}

If a code assigns lengths \(l_i\) to symbol \(i\), the average length is:
\begin{equation}
L = \sum_{i=1}^{N} p_i \, l_i
\end{equation}

An optimal code satisfies:
\[
H \leq L < H + 1
\]

Huffman and arithmetic coding approach this bound.

\subsection{Prefix Codes}

A prefix code is a code where no codeword is a prefix of another.

\textbf{Properties:}
\begin{itemize}
\item Instantaneously decodable  
\item Satisfy Kraft-McMillan inequality:
\[
\sum_{i=1}^{N} 2^{-l_i} \leq 1
\]
\item Basis for Huffman coding and many entropy codes
\end{itemize}

\subsection{Lossless vs. Lossy Compression}

\subsubsection{Lossless Compression}

\begin{itemize}
\item Exact reconstruction
\item Used for PNG, GIF, medical imaging, LZ77/LZW, text files
\end{itemize}

\subsubsection{Lossy Compression}

\begin{itemize}
\item Allows controlled distortion
\item Used in JPEG, MPEG
\item Removes psychovisual redundancy
\end{itemize}

\subsection{Modeling Data for Compression}

There are two broad approaches:

\begin{enumerate}
\item \textbf{Statistical modeling} — probabilities assigned to symbols  
Used in Huffman, arithmetic coding.

\item \textbf{Dictionary modeling} — recurring patterns replaced with references  
Used in LZ77, LZ78, LZW, PNG.
\end{enumerate}

Both approaches appear in image compression:

\begin{itemize}
\item JPEG uses statistical + transform coding.
\item PNG uses LZ77.
\item TIFF can use LZW.
\end{itemize}

\subsection{Two-Part Structure of Compression Algorithms}

Most compression systems consist of:

\begin{enumerate}
\item \textbf{A redundancy reduction step}  
(e.g., transform coding, prediction, dictionary encoding)

\item \textbf{An entropy coding step}  
(Huffman or arithmetic coding)
\end{enumerate}

For example, JPEG pipeline:

\[
\text{DCT} \rightarrow \text{Quantization} \rightarrow \text{Zig-zag} \rightarrow \text{Huffman}
\]

\subsection{Example Problem 1: Compute Entropy and Minimum Bits Required}

\subsubsection*{Problem}

A source emits symbols with probabilities:
\[
p = \{0.5,\, 0.25,\, 0.125,\, 0.125\}
\]

Compute:
\begin{enumerate}
\item Entropy
\item Minimum number of bits per symbol (lower bound)
\item Redundancy if average codeword length is 2 bits
\end{enumerate}

\subsubsection*{Solution}

\paragraph{1. Entropy}
\[
H = -\sum p_i\log_2 p_i
\]

\[
= -[0.5(-1) + 0.25(-2) + 0.125(-3) + 0.125(-3)]
\]

\[
= 0.5 + 0.5 + 0.375 + 0.375 = 1.75 \text{ bits}
\]

\paragraph{2. Minimum bits:}
\[
L_{\min} = H = 1.75
\]

\paragraph{3. Redundancy:}
\[
R = L - H = 2 - 1.75 = 0.25 \text{ bits/symbol}
\]

\subsubsection*{Interpretation}

Even the best code cannot do better than 1.75 bits/symbol.

\subsection{Example Problem 2: Verify Kraft Inequality and Prefix Condition}

\subsubsection*{Problem}

Consider the following code:

\[
A: 0,\quad
B: 10,\quad
C: 110,\quad
D: 111
\]

\begin{enumerate}
\item Verify Kraft inequality.
\item Determine if this is a prefix-free code.
\end{enumerate}

\subsubsection*{Solution}

\paragraph{1. Kraft inequality}

\[
K = 2^{-1} + 2^{-2} + 2^{-3} + 2^{-3}
\]

\[
= 0.5 + 0.25 + 0.125 + 0.125 = 1.0
\]

Since \(K \le 1\), the code is valid.

\paragraph{2. Prefix condition}

Check if any codeword is prefix of another:

\begin{itemize}
\item 0 is not a prefix of any word starting with 1.
\item 10 is not a prefix of 110 or 111.
\item 110 and 111 do not overlap.
\end{itemize}

Thus, the code is \textbf{prefix-free}.

\subsubsection*{Interpretation}

The code is uniquely decodable and instantaneous.

\subsection{Summary}

\begin{itemize}
\item Compression reduces redundancy (coding, interpixel, psychovisual).
\item Entropy defines the theoretical lower bound for lossless compression.
\item Prefix-free codes enable instant decoding.
\item Compression uses modeling (dictionary/statistical) + entropy coding.
\item This foundation underlies LZ77, LZ78, LZW, Huffman, arithmetic, and Golomb-Rice coding.
\end{itemize}

\section{LZ77 Compression}

\subsection{Introduction}

LZ77 is a dictionary-based, lossless compression algorithm introduced by Lempel and Ziv in 1977.  
It forms the foundation of many modern compressors, including:

\begin{itemize}
\item \textbf{DEFLATE} (used in PNG, ZIP, gzip)
\item \textbf{LZSS}, \textbf{LZMA}, \textbf{Zstd}, \textbf{LZO}
\end{itemize}

LZ77 replaces repeated sequences with \textbf{pointers} referencing previous occurrences within a sliding window.

\subsection{Sliding Window Model}

LZ77 uses:
\begin{itemize}
\item a \textbf{search buffer}: previously encoded text
\item a \textbf{look-ahead buffer}: text to be encoded next
\end{itemize}

\[
\underbrace{\text{past data}}_{\text{search buffer}}
\;|\;
\underbrace{\text{future data}}_{\text{look-ahead buffer}}
\]

The encoder outputs triples of the form:

\[
(d, \ell, s)
\]

Where:
\begin{itemize}
\item \(d\) = distance backward (offset)
\item \(\ell\) = length of match
\item \(s\) = next unmatched symbol
\end{itemize}

They mean:  
“Go back \(d\) characters, copy \(\ell\) characters, then add symbol \(s\).”

If no match is found:
\[
(0, 0, s)
\]
is emitted.

\subsection{Matching Procedure}

\begin{enumerate}
\item For the start of the lookahead buffer, search for the longest match in the search buffer.
\item Compute distance \(d\) and length \(\ell\).
\item Emit triple.
\item Slide window by \(\ell + 1\) positions.
\end{enumerate}

\subsection{Decoding}

Decoding simply reverses the pointer:

\[
\text{If } (d,\ell,s), \text{ then output:}
\underbrace{\text{copy previous }\ell\text{ characters from offset } d}_{\text{dictionary copy}}, 
\text{ then write symbol } s.
\]

\subsection{Why LZ77 Works}

Natural language text and image data (after filtering/transforming) contain repeated patterns. Examples:

\begin{itemize}
\item Repeated words, prefixes, suffixes.
\item Repeated byte sequences in PNG filtered rows.
\item Runs of similar symbols after prediction.
\end{itemize}

Replacing these with backward references achieves significant compression.

\subsection{LZ77 in Image Processing}

PNG uses DEFLATE = (LZ77 + Huffman coding).  
Before LZ77, PNG applies a reversible line predictor to convert pixel values into differences → increasing redundancy.

Thus LZ77 is indirectly important for lossless image compression.

\subsection{Example Problem 1: Manual LZ77 Encoding}

\subsubsection*{Problem}

Encode the string:
\[
\text{BANANA BAN}
\]

Assume unlimited window size.

\subsubsection*{Solution}

Start:

\textbf{B}: no match  
\[
(0,0,B)
\]

\textbf{A}: no match  
\[
(0,0,A)
\]

\textbf{N}: no match  
\[
(0,0,N)
\]

\textbf{A}: previous A at position 2 → match length 1
\[
d = 3,\;\ell = 1,\; s = N
\]
So:
\[
(3,1,N)
\]

Next sequence: “ANA BAN”

Match “ANA” with earlier “ANA” starting at position 2:

\[
d=4,\;\ell=3,\;s=\ \text{space}
\]

Output:
\[
(4,3,\text{space})
\]

Remaining: “BAN”

We saw “BAN” at the beginning.

\[
d=9,\;\ell=3,\;s=\text{EOF}
\]

\subsubsection*{Result}

\[
(0,0,B),\; (0,0,A),\; (0,0,N),\; (3,1,N),\; (4,3,\text{space}),\; (9,3,\text{EOF})
\]

\subsubsection*{Interpretation}

LZ77 replaces repeated substrings with compact references.

\subsection{Example Problem 2: LZ77 Decoding}

\subsubsection*{Problem}

Decode the LZ77 output:

\[
(0,0,h),\;
(0,0,e),\;
(2,1,l),\;
(3,2,o)
\]

\subsubsection*{Solution}

Step-by-step:

\paragraph{1.} \((0,0,h)\) → output "h"

\paragraph{2.} \((0,0,e)\) → "he"

\paragraph{3.} \((2,1,l)\)  
Look back 2 characters: "he"  
Copy 1 → "h"  
Output:
\[
\text{"hel"}
\]

\paragraph{4.} \((3,2,o)\)  
Look back 3 characters: substring "hel"  
Copy 2: "he"  
Append "o":  
\[
\text{hello}
\]

So final output:

\[
\boxed{\text{hello}}
\]

\subsection{Summary}

\begin{itemize}
\item LZ77 uses a sliding window to replace repeated sequences with backward references.
\item Output triples: \((d,\ell,s)\)
\item Basis of PNG, ZIP, gzip.
\item Very effective after predictive filtering (as done in image compression).
\item Decoding is extremely simple: copy from history and append next symbol.
\end{itemize}

\section{LZ78 Compression}

\subsection{Introduction}

LZ78 (Lempel–Ziv, 1978) is a lossless dictionary-based compression algorithm.  
Unlike LZ77, which uses a sliding window, LZ78 builds an explicit dictionary of previously seen \emph{phrases}.  
Each new phrase is assigned an index and added to the dictionary.

LZ78 is fundamental in compression theory and is the basis of:

\begin{itemize}
\item LZW (used in GIF and some TIFF variants)
\item Many adaptive dictionary compressors
\end{itemize}

\subsection{Encoding Mechanism}

LZ78 builds a dictionary of the form:

\[
1:\text{phrase}_1,\quad 2:\text{phrase}_2,\quad \ldots
\]

Encoding outputs \textbf{pairs}:

\[
(i, c)
\]

Where:
\begin{itemize}
\item \(i\) = dictionary index of the longest matching phrase  
\item \(c\) = next symbol after that phrase
\end{itemize}

Meaning:
\[
\text{dictionary}[i] + c
\]
forms the next new phrase.

\subsubsection{Encoding Steps}

Given input string \(S\):

\begin{enumerate}
\item Initialize dictionary with index 0 representing the empty string.
\item Scan input for the longest phrase present in dictionary.
\item Let match index be \(i\), next character be \(c\).
\item Output pair: \((i, c)\).
\item Insert new phrase \( \text{dictionary}[i] + c \) into dictionary.
\item Move forward by the length of the new phrase.
\end{enumerate}

\subsection{Decoding}

Decoding uses the same dictionary-building process, but reconstructs phrases as follows:

Given pair \((i, c)\):

\[
\text{phrase} = \text{dictionary}[i] + c
\]
Append this phrase to the output, then add it to the dictionary.

\subsection{Comparison: LZ77 vs. LZ78}

\begin{table}[h!]
\centering
\begin{tabular}{|c|c|}
\hline
\textbf{LZ77} & \textbf{LZ78} \\
\hline
Sliding window dictionary & Explicit, growing dictionary \\
\hline
Dictionary entries overlap & Dictionary entries are stored explicitly \\
\hline
Output: (distance, length, next) & Output: (index, next) \\
\hline
Used by PNG (DEFLATE) & Basis for LZW (GIF, TIFF) \\
\hline
Better for repeated local patterns & Better for repetitive global patterns \\
\hline
\end{tabular}
\end{table}

\subsection{Example Problem 1: LZ78 Encoding}

\subsubsection*{Problem}

Encode the string:

\[
\text{ABABABA}
\]

using LZ78.

\subsubsection*{Solution}

Start with dictionary:

\[
0: \epsilon
\]

\textbf{Step 1:}  
Longest match = empty string, next char = A  
Output:
\[
(0, A)
\]
Add phrase 1:
\[
1: A
\]

\textbf{Step 2:}  
Next input starts with B  
Match = empty, next = B  
\[
(0, B)
\]
Add phrase 2:
\[
2: B
\]

\textbf{Step 3:}  
Next input: ABABA  
Match = A (index 1), next = B  
\[
(1, B)
\]
Add phrase 3:
\[
3: AB
\]

\textbf{Step 4:}  
Remaining: ABA  
Match = AB (index 3), next = A  
\[
(3, A)
\]
Add phrase 4:
\[
4: ABA
\]

Encoding complete.

\subsubsection*{Final Output}

\[
(0,A),\; (0,B),\; (1,B),\; (3,A)
\]

\subsubsection*{Interpretation}

The dictionary stores meaningful repeated phrases:
\[
A,\; B,\; AB,\; ABA
\]

\subsection{Example Problem 2: LZ78 Decoding}

\subsubsection*{Problem}

Decode the LZ78 sequence:

\[
(0, A),\; (0, B),\; (1, C),\; (3, B)
\]

\subsubsection*{Solution}

Initialize:
\[
0:\epsilon
\]

\textbf{1. } (0, A)  
Phrase = \(\epsilon + A = A\)  
Output: A  
Add to dictionary:
\[
1: A
\]

\textbf{2. } (0, B)  
Phrase = \(\epsilon + B = B\)  
Output: B  
Add:
\[
2: B
\]

\textbf{3. } (1, C)  
Dictionary[1] = A  
Phrase = AC  
Output: AC  
Add:
\[
3: AC
\]

\textbf{4. } (3, B)  
Dictionary[3] = AC  
Phrase = ACB  
Output: ACB  
Add:
\[
4: ACB
\]

\subsubsection*{Final Output}

\[
\boxed{A\, B\, AC\, ACB}
\]

Concatenating:

\[
\boxed{\text{ABACACB}}
\]

\subsection{Summary}

\begin{itemize}
\item LZ78 uses explicit dictionary entries indexed by integers.
\item Output consists of \((\text{index}, \text{next symbol})\) pairs.
\item Dictionary grows as new phrases are discovered.
\item More efficient than LZ77 when repeated phrases appear across long distances.
\item Basis of the widely used LZW compression algorithm.
\end{itemize}

\section{LZW Compression}

\subsection{Introduction}

LZW (Lempel–Ziv–Welch, 1984) is an improved variant of LZ78.  
It is a \textbf{lossless}, dictionary-based compression algorithm widely used in:

\begin{itemize}
\item GIF images
\item TIFF formats
\item UNIX compress utility
\item PDF streams (optional)
\end{itemize}

LZW modifies LZ78 so that the encoder outputs \textbf{only index codes}, not (index, symbol) pairs.  
It relies on the fact that the decoder can reconstruct the same dictionary dynamically.

\subsection{Key Idea}

LZW initializes the dictionary with all possible single-character symbols.  
For 8-bit data (e.g., grayscale or GIF-style indexed color):

\[
\text{Dictionary initially contains entries 0–255}
\]

As encoding proceeds, longer strings are added.

\subsection{Encoding Algorithm}

Given input string \(S\):

\begin{enumerate}
\item Initialize dictionary with all single characters (256 entries).
\item Set \(w\) = first character of input.
\item Read next character \(k\).
\item If \(wk\) exists in dictionary, set \(w \leftarrow wk\).
\item Else:
	\begin{itemize}
	\item Output code for \(w\)
	\item Add \(wk\) to dictionary
	\item Set \(w = k\)
	\end{itemize}
\item Repeat until end of input, then output code for last \(w\).
\end{enumerate}

\subsection{Decoding Algorithm}

Given code stream:

\begin{enumerate}
\item Initialize dictionary with 256 single-character entries.
\item Read first code → output its character → set \(w\) to that character.
\item For each subsequent code \(k\):
	\begin{itemize}
	\item If \(k\) is in dictionary: entry = dictionary[k]
	\item Else (edge case): entry = \(w + w[0]\)
	\item Output entry
	\item Add \(w + entry[0]\) to dictionary
	\item Set \(w = entry\)
	\end{itemize}
\end{enumerate}

\subsection{Why LZW Works Well}

\begin{itemize}
\item Single-character dictionary initialization ensures fast start.
\item Long repeated patterns (common in text and image scanlines) are encoded as single codes.
\item Dictionary entries grow to reflect statistical structure of data.
\item Efficient for images with repeated colors, lines, patterns (e.g., cartoons, diagrams).
\end{itemize}

\subsection{LZW in Image Processing (GIF/TIFF)}

GIF operates on indexed images (0–255 color indices).  
Process:

\begin{enumerate}
\item Each pixel value is a symbol in \([0,255]\).
\item LZW compresses the pixel stream.
\item No loss of image fidelity.
\end{enumerate}

TIFF offers LZW as an optional lossless compression mode.

\subsection{Example Problem 1: LZW Encoding}

\subsubsection*{Problem}

Encode the string:

\[
\text{ABABABA\$}
\]

Assume ASCII values 65=A, 66=B, 36=\$.

\subsubsection*{Solution}

Initial dictionary: 0–255.

Start:

\textbf{w = A}  
Next symbol: B  
"A B" not in dictionary → output A=65  
Add "AB" at index 256  
Set w = B

\textbf{w = B}  
Next: A  
"BA" not in dictionary → output 66  
Add "BA" at 257  
w = A

\textbf{w = A}  
Next: B  
"AB" is in dictionary (256)  
So set w = "AB"

Next: A  
"ABA" not in dictionary → output code for "AB" = 256  
Add "ABA" at index 258  
Set w = A

Next: \$  
"A\$" not in dictionary → output A=65  
Add "A\$" at 259  
Set w = \$

End of input → output \$ = 36

\subsubsection*{Final Encoded Output}

\[
\boxed{65,\; 66,\; 256,\; 65,\; 36}
\]

\subsubsection*{Interpretation}

Repeated patterns "AB" and "ABA" were detected automatically and encoded as single dictionary references.

\subsection{Example Problem 2: LZW Decoding}

\subsubsection*{Problem}

Decode the LZW stream:

\[
65,\; 66,\; 256,\; 258,\; 36
\]

Assume same dictionary initialization (ASCII).

\subsubsection*{Solution}

Initial dictionary: entries 0–255.

\textbf{Step 1:}  
Code 65 → "A"  
Output: A  
Set w = "A"

\textbf{Step 2:}  
Code 66 → "B"  
Output: B  
Add "AB" (w + B[0]) at 256  
w = "B"

\textbf{Step 3:}  
Code 256 → dictionary[256] = "AB"  
Output: AB  
Add "BA" at 257 (w + A)  
w = "AB"

\textbf{Step 4:}  
Code 258 → not yet defined.  
Special case: entry = w + w[0] = "AB" + "A" = "ABA"  
Output: ABA  
Add "AB" + "A" = "ABA" at 258  
w = "ABA"

\textbf{Step 5:}  
Code 36 → "\$"  
Output: \$  
Add "ABA" + "\$" at 259  

\subsubsection*{Final Output}

Concatenating:

\[
\boxed{\text{ABABABA\$}}
\]

\subsection{Summary}

\begin{itemize}
\item LZW initializes dictionary with all symbols, unlike LZ78.
\item Output is a sequence of indices only.
\item Dictionary is built identically during encoding and decoding.
\item Widely used in GIF and TIFF image formats.
\item Efficient for repeated raster patterns in images.
\end{itemize}

\section{Encoding Theory}

\subsection{Introduction}

Encoding is the foundation of all lossless compression systems.  
Given a set of symbols \( \mathcal{A} = \{a_1, a_2, \ldots, a_n\} \) with associated probabilities
\[
P = \{p_1, p_2, \ldots, p_n\},
\]
the goal is to assign binary codewords
\[
C = \{c_1, c_2, \ldots, c_n\}
\]
such that the expected number of bits per symbol is minimized.

\[
L = \sum_{i=1}^{n} p_i \, \ell_i
\]
where \( \ell_i \) is the length of codeword \(c_i\).

\subsection{Uniquely Decodable and Prefix Codes}

A code is \textbf{uniquely decodable} if every encoded bitstream has exactly one valid decoding.

A more restrictive class, \textbf{prefix codes}, ensures:

\[
\text{No codeword is a prefix of any other codeword.}
\]

Prefix codes (e.g., Huffman codes) allow instantaneous decoding:
\begin{itemize}
\item No look-ahead required
\item No ambiguity
\item Decoding can proceed bit by bit
\end{itemize}

Every prefix code corresponds to a binary tree:
\begin{itemize}
\item Leaves = symbols
\item Path to leaf = codeword
\item Left child = 0, right child = 1 (or vice versa)
\end{itemize}

\subsection{Kraft–McMillan Inequality}

For any prefix code with codeword lengths \(\ell_1, \ell_2, \ldots, \ell_n\):

\[
\boxed{\sum_{i=1}^{n} 2^{-\ell_i} \le 1}
\]

Conversely, if a set of lengths satisfies Kraft’s inequality, a prefix code exists.

\paragraph{Proof Sketch:}
In a full binary tree, a node at depth \(\ell\) occupies a fraction \(2^{-\ell}\) of the tree capacity.  
All codewords must fit into the tree without overlap → total “capacity” used must not exceed 1.

This inequality is the basis for validating or designing prefix codes.

\subsection{Optimality of Prefix Codes}

Shannon’s source coding theorem states:

\[
\boxed{H(X) \le L < H(X) + 1}
\]
where
\[
H(X) = -\sum p_i \log_2 p_i
\]
is the entropy of the source.

A code is optimal when it minimizes expected code length \(L\).  
Huffman coding achieves the optimal prefix code.

\subsection{Entropy and Average Code Length}

\[
H(X) = -\sum_{i=1}^{n} p_i \log_2 p_i
\]

The lower bound for any uniquely decodable binary code is:

\[
L \ge H(X)
\]

If probabilities are powers of \(1/2\), Huffman coding achieves:

\[
L = H(X)
\]

\subsection{Instantaneous Codes}

Instantaneous codes are exactly the prefix codes.

Properties:
\begin{itemize}
\item Decoded without delay
\item No ambiguity
\item Represented by full binary trees
\end{itemize}

Example:
\[
a = 0,\; b = 10,\; c = 110,\; d = 111
\]

No codeword is the prefix of any other → prefix-free.

\subsection{Canonical Codes}

Some compressors (e.g., DEFLATE) use \textbf{canonical Huffman codes}, which:
\begin{itemize}
\item Assign shorter codes to smaller indices
\item Require only code lengths to reconstruct full table
\item Provide deterministic compact representation
\end{itemize}

Canonical Huffman is vital for PNG, ZIP, gzip.

\subsection{Example Problem 1: Checking Prefix Validity}

\subsubsection*{Problem}

Determine whether the following code is a prefix code:

\[
a=0,\quad b=01,\quad c=011,\quad d=111
\]

\subsubsection*{Solution}

Check if any codeword is a prefix of another.

\begin{itemize}
\item \(a = 0\) is a prefix of all others (\(01, 011\)).
\end{itemize}

Thus this is \textbf{not} a prefix code.

Additionally, Kraft inequality check:

\[
2^{-1} + 2^{-2} + 2^{-3} + 2^{-3} = \frac{1}{2} + \frac{1}{4} + \frac{1}{8} + \frac{1}{8} = 1
\]

Even though Kraft equality holds, prefix condition fails → code is uniquely decodable but not instantaneous.

\subsubsection*{Conclusion}

\[
\boxed{\text{Not a prefix code.}}
\]

\subsection{Example Problem 2: Designing a Prefix Code From Given Lengths}

\subsubsection*{Problem}

Given symbol set \( \{A, B, C, D\} \) with desired codeword lengths:

\[
\ell_A = 1,\; \ell_B = 2,\; \ell_C = 3,\; \ell_D = 3
\]

(a) Check if a prefix code is possible.  
(b) Construct a valid prefix code.

\subsubsection*{Solution}

\paragraph{(a) Kraft Check}

\[
2^{-1} + 2^{-2} + 2^{-3} + 2^{-3}
= \frac{1}{2} + \frac{1}{4} + \frac{1}{8} + \frac{1}{8}
= 1
\]

Kraft inequality satisfied → prefix code exists.

\paragraph{(b) Constructing Code}

Assign shortest codes to earliest leaves:

\[
A = 0
\]

Remaining tree capacity starts with prefixes beginning with “1”.

\[
B = 10
\]

Next, two leaves at depth 3:

\[
C = 110,\quad D = 111
\]

\subsubsection*{Final Prefix Code}

\[
\boxed{A=0,\; B=10,\; C=110,\; D=111}
\]

\subsection{Summary}

\begin{itemize}
\item Encoding assigns binary codewords to symbols to minimize expected length.
\item Prefix codes ensure instantaneous decoding.
\item Kraft inequality validates prefix code feasibility.
\item Shannon entropy gives the lower bound on average code length.
\item Canonical codes allow compact representation and efficient decoding.
\end{itemize}
\section{Huffman Coding}

\subsection{Introduction}

Huffman Coding (David Huffman, 1952) is a \textbf{lossless}, \textbf{optimal} prefix-code algorithm for symbols with known probabilities.  
It generates a binary tree where:

\begin{itemize}
\item high-probability symbols receive short codewords,
\item low-probability symbols receive long codewords.
\end{itemize}

Huffman coding minimizes the expected code length:
\[
L = \sum_{i} p_i \ell_i
\]
subject to the prefix condition.

It is widely used in:

\begin{itemize}
\item JPEG (as entropy coding after DCT quantization)
\item PNG/ZIP/gzip (DEFLATE uses \textit{canonical} Huffman)
\item TIFF
\end{itemize}

\subsection{Constructing the Huffman Code}

Given symbols \(a_1, a_2, \ldots, a_n\) with probabilities \(p_1, p_2, \ldots, p_n\):

\begin{enumerate}
\item Create a leaf node for each symbol with weight \(p_i\).
\item Insert all nodes into a min-priority queue.
\item Repeat until a single node remains:
	\begin{itemize}
	\item Remove two nodes with smallest weights.
	\item Create new node with weight = sum of weights.
	\item Attach the removed nodes as children.
	\item Insert the new node back into the queue.
	\end{itemize}
\item Assign:
	\[
	\text{left child} \rightarrow 0,\qquad \text{right child} \rightarrow 1
	\]
\item Read codes from root to leaves.
\end{enumerate}

\subsection{Properties of Huffman Codes}

\paragraph{Optimality}
Huffman codes minimize expected length among all prefix codes:
\[
L_{\text{Huffman}} = \min L
\]

\paragraph{Prefix-free}
No codeword is a prefix of any other.

\paragraph{Full Binary Tree}
Every internal node has exactly two children.

\paragraph{Length Bounds}
\[
H(X) \le L < H(X) + 1
\]

\paragraph{Limitations}
Huffman coding is suboptimal when:
\begin{itemize}
\item probabilities change over time (adaptive schemes needed),
\item symbol frequencies are uniform (little gain),
\item probabilities are extremely skewed (arithmetic coding may outperform).
\end{itemize}

\subsection{Canonical Huffman Coding (Used in PNG/ZIP)}

Canonical Huffman does not store tree structure.  
Instead it stores only \emph{code lengths}.  

Decoder reconstructs codes using rules:
\begin{enumerate}
\item Sort symbols by increasing length.
\item Shortest codes start from all zeros.
\item Each subsequent code is binary increment of previous.
\end{enumerate}

Canonical coding ensures:
\begin{itemize}
\item deterministic assignment,
\item compact representation,
\item faster decoding via table lookup.
\end{itemize}

\subsection{Example Problem 1: Full Huffman Construction}

\subsubsection*{Problem}

Given symbols and probabilities:

\[
A:0.40,\quad B:0.20,\quad C:0.20,\quad D:0.10,\quad E:0.10
\]

Construct the Huffman code and compute average length.

\subsubsection*{Solution}

\paragraph{Step 1: Combine smallest probabilities}

D (0.10) + E (0.10) → DE (0.20)

Now we have:

\[
A:0.40,\; B:0.20,\; C:0.20,\; DE:0.20
\]

\paragraph{Step 2: Combine next smallest}

Choose any two of 0.20:

\[
B(0.20) + C(0.20) \rightarrow BC(0.40)
\]

\paragraph{Step 3: Combine remaining lowest}

\[
DE(0.20) + (BC)(0.40) \rightarrow DEBC(0.60)
\]

\paragraph{Step 4: Combine last two}

\[
A(0.40) + (DEBC)(0.60) \rightarrow 1.00
\]

\paragraph{Assign bits}

Choose:

Left child = 0  
Right child = 1  

Resulting codes (one valid assignment):

\[
\begin{aligned}
A &= 0 \\
B &= 100 \\
C &= 101 \\
D &= 1100 \\
E &= 1101
\end{aligned}
\]

\paragraph{Compute Average Length}

\[
L = 0.40(1) + 0.20(3) + 0.20(3) + 0.10(4) + 0.10(4)
\]

\[
L = 0.4 + 0.6 + 0.6 + 0.4 + 0.4 = 2.4
\]

\subsubsection*{Final Answer}

\[
\boxed{
A=0,\; B=100,\; C=101,\; D=1100,\; E=1101,\; L=2.4
}
\]

\subsection{Example Problem 2: Canonical Huffman Code Construction}

\subsubsection*{Problem}

Symbols have the following codeword lengths:

\[
\ell_A=2,\; \ell_B=2,\; \ell_C=3,\; \ell_D=3,\; \ell_E=3
\]

Construct the canonical Huffman code.

\subsubsection*{Solution}

Step 1: Sort symbols by length (and alphabetically if tied):

\[
A(2), B(2), C(3), D(3), E(3)
\]

Step 2: Assign smallest code:

\[
A = 00
\]

Step 3: Next code is binary increment:

\[
B = 01
\]

Step 4: Increase length to 3 bits:

Append trailing zeros when moving to next length:

\[
C = 100,\quad D = 101,\quad E = 110
\]

\subsubsection*{Final Canonical Codes}

\[
\boxed{
A=00,\; B=01,\; C=100,\; D=101,\; E=110
}
\]

\subsection{Summary}

\begin{itemize}
\item Huffman coding produces optimal prefix codes.
\item Average code length satisfies \(H \le L < H+1\).
\item Construction uses bottom-up merging of least probable symbols.
\item Canonical Huffman codes allow efficient storage and decoding.
\item Widely used in image compression (JPEG / PNG).
\end{itemize}

\section{Arithmetic Coding}

\subsection{Introduction}

Arithmetic coding is a form of entropy coding that represents an entire message as a single number in the interval
\([0, 1)\).  
Unlike Huffman coding, which assigns a fixed number of bits per symbol, arithmetic coding assigns \emph{fractional} bits by encoding the entire sequence into a progressively refined interval.

It is widely used in:
\begin{itemize}
\item JPEG2000
\item H.264 / H.265 video coding (CABAC)
\item Modern image codecs such as AVIF, WebP, HEIC
\end{itemize}

Arithmetic coding achieves compression rates extremely close to the entropy limit.

\subsection{Basic Concept}

Given a source alphabet:
\[
\mathcal{A} = \{a_1, a_2, \ldots, a_n\}
\]
with probabilities:
\[
P = \{p_1, p_2, \ldots, p_n\}
\]

We construct the cumulative distribution:

\[
C_0 = 0,\quad C_i = \sum_{k=1}^{i} p_k
\]

Each symbol \(a_i\) corresponds to the subinterval:
\[
[C_{i-1},\, C_i)
\]

\subsection{Encoding Procedure}

Arithmetic coding maintains a narrowing interval:
\[
[\,L,\, U\,)
\]
initially:
\[
L=0,\quad U=1.
\]

For each symbol \(s\):

\begin{enumerate}
\item Compute current range:
\[
R = U - L
\]
\item Set:
\[
U \leftarrow L + R \cdot C_{s,\text{high}}
\]
\[
L \leftarrow L + R \cdot C_{s,\text{low}}
\]
\end{enumerate}

After encoding all symbols, any number in the final interval uniquely represents the message.

\subsection{Decoding Procedure}

Decoder receives a number \(x \in [L, U)\).  
At each step:

\begin{enumerate}
\item Determine which subinterval contains \(x\).
\item Output the corresponding symbol.
\item Update interval using same formulas as encoder.
\item Repeat until message length is reached.
\end{enumerate}

\subsection{Advantages Over Huffman Coding}

\begin{itemize}
\item \textbf{No requirement for symbol code lengths to be integers}.
\item \textbf{Compression efficiency near entropy}:
\[
L \approx H(X)
\]
\item Easily extends to:
	\begin{itemize}
	\item adaptive models
	\item context-based probability estimation
	\item high-order Markov models
	\end{itemize}
\item Crucial for modern codecs handling complex symbol statistics.
\end{itemize}

\subsection{Renormalization (Practical Implementation)}

In practice, intervals become extremely small.  
To avoid underflow:

\begin{itemize}
\item Range is rescaled when \(L\) and \(U\) fall into the same half or quarter.
\item Bits are emitted progressively.
\item This is done using integer arithmetic.
\end{itemize}

This is the basis of:
\begin{itemize}
\item MQ-coder (JPEG2000)
\item CABAC (H.264/AVC, H.265/HEVC)
\end{itemize}

\subsection{Example Problem 1: Arithmetic Encoding}

\subsubsection*{Problem}

Symbols and probabilities:

\[
A:0.5,\quad B:0.3,\quad C:0.2
\]

Cumulative intervals:

\[
A:[0,0.5),\; B:[0.5,0.8),\; C:[0.8,1.0)
\]

Encode the message:

\[
\text{``ABC''}
\]

\subsubsection*{Solution}

\paragraph{Initial:} \(L=0,\; U=1\)

\textbf{Step 1: Encode A}

\[
R = 1-0 = 1
\]
\[
L = 0 + 1\cdot 0 = 0,\quad U = 0 + 1\cdot 0.5 = 0.5
\]

\textbf{New interval:} \([0, 0.5)\)

\textbf{Step 2: Encode B}

\[
R = 0.5 - 0 = 0.5
\]
\[
L = 0 + 0.5\cdot 0.5 = 0.25
\]
\[
U = 0 + 0.5\cdot 0.8 = 0.40
\]

\textbf{New interval:} \([0.25, 0.40)\)

\textbf{Step 3: Encode C}

\[
R = 0.40 - 0.25 = 0.15
\]
\[
L = 0.25 + 0.15\cdot 0.8 = 0.37
\]
\[
U = 0.25 + 0.15\cdot 1.0 = 0.40
\]

Final interval:
\[
[0.37,\, 0.40)
\]

Any number in this interval—commonly the lower bound:

\[
\boxed{0.37}
\]

represents the full message “ABC”.

\subsection{Example Problem 2: Arithmetic Decoding}

\subsubsection*{Problem}

Use the same probability model as above.  
Decode the number:

\[
x = 0.381
\]

Message length = 3 symbols.

\subsubsection*{Solution}

\paragraph{Initial interval:} \([0,1)\)

\textbf{Step 1: Find symbol for } x = 0.381

Intervals:
\[
A:[0,0.5),\quad B:[0.5,0.8),\quad C:[0.8,1)
\]

Since \(0.381 \in A\), output:
\[
A
\]

Update interval to \([0,0.5)\).

\textbf{Step 2: Decode within interval \([0,0.5)\)}

Normalize:
\[
x' = \frac{x - 0}{0.5 - 0} = 0.381 / 0.5 = 0.762
\]

Check subintervals of A/B/C:

\[
A:[0,0.5),\quad B:[0.5,0.8),\quad C:[0.8,1)
\]

Since \(0.762 \in B\), symbol = B.

Update interval to:
\[
[0.25, 0.40)
\]

\textbf{Step 3: Decode last symbol}

Normalize:
\[
x'' = \frac{0.381 - 0.25}{0.40 - 0.25} = \frac{0.131}{0.15} \approx 0.873
\]

\(0.873 \in C\) interval.

\textbf{Output: C}

\subsubsection*{Final Decoded Message}

\[
\boxed{\text{ABC}}
\]

\subsection{Summary}

\begin{itemize}
\item Arithmetic coding represents a sequence as a progressively shrinking interval.
\item Achieves compression near the entropy limit.
\item More powerful than Huffman coding for non-integer probabilities.
\item Forms the basis of modern image/video compression systems.
\item Encoding = refining \([L,U)\) interval; decoding = locating \(x\) within nested intervals.
\end{itemize}

\section{Golomb--Rice Coding}

\subsection{Introduction}

Golomb coding is a lossless entropy coding method designed for sources whose symbol distribution resembles a \textbf{geometric distribution}:

\[
P(X = k) = (1-p)^k p, \qquad k = 0,1,2,\ldots
\]

Golomb codes are optimal for such distributions and are widely used in:

\begin{itemize}
\item Lossless image compression (JPEG-LS)
\item Audio compression (FLAC residual coding)
\item Predictive coding of images (prediction errors often geometric/Laplacian)
\end{itemize}

A special case of Golomb coding, called \textbf{Rice coding}, occurs when the divisor parameter is a power of two.

\subsection{Golomb Coding Overview}

Golomb coding uses a parameter \(m\), chosen based on the statistical distribution of the data.  
For a non-negative integer \(N\), compute:

\[
q = \left\lfloor \frac{N}{m} \right\rfloor, \qquad r = N \bmod m
\]

The Golomb code for \(N\) is:

\[
\text{code}(N) = \underbrace{11\ldots1}_{q \text{ times}}0 \; \| \; \text{binary\_code}(r)
\]

The first part is unary coding of \(q\), followed by a truncated binary encoding of the remainder \(r\).

\subsection{Truncated Binary Encoding of the Remainder}

Let:
\[
b = \lceil \log_2 m \rceil
\]
and:
\[
t = 2^b - m
\]

Then encoding of remainder \(r\) is:

\[
\text{binary\_code}(r) = 
\begin{cases}
\text{binary}(r), & r < t \\
\text{binary}(r + t), & r \ge t
\end{cases}
\]

\subsection{Decoding}

Given a Golomb code:

\begin{enumerate}
\item Count number of leading 1s until the first 0 → this gives \(q\).
\item Decode the remaining bits to obtain \(r\) (undo truncated binary).
\item Reconstruct value:
\[
N = qm + r
\]
\end{enumerate}

\subsection{Choosing the Parameter \(m\)}

For geometric distributions:
\[
m \approx -\frac{1}{\ln(1-p)}
\]

In image compression (JPEG-LS), the optimal \(m\) is computed adaptively from the variance of prediction errors.

\subsection{Rice Coding (Special Case)}

Rice coding is a simplified Golomb code where:

\[
m = 2^k
\]

Thus:

\[
q = N >> k \quad (\text{right shift})
\]
\[
r = N \& (2^k - 1)
\]

Binary representation of \(r\) is simply the \(k\)-bit binary number.

This makes Rice coding extremely efficient for hardware and software implementations.

\subsection{Example Problem 1: Golomb Encoding}

\subsubsection*{Problem}

Encode the integer \(N = 37\) using Golomb coding with parameter \(m = 10\).

\subsubsection*{Solution}

Compute:

\[
q = \left\lfloor \frac{37}{10} \right\rfloor = 3,\qquad r = 37 \bmod 10 = 7
\]

Unary code for \(q=3\):

\[
1110
\]

Now encode remainder:

\[
b = \lceil \log_2 10 \rceil = 4,\quad t = 2^4 - 10 = 6
\]

Since \(r = 7 \ge t = 6\), encode:
\[
r + t = 7 + 6 = 13
\]
Binary(13) in 4 bits:
\[
1101
\]

\subsubsection*{Final Golomb Code}

\[
\boxed{1110\,1101}
\]

\subsection{Example Problem 2: Rice Coding}

\subsubsection*{Problem}

Encode the integer sequence using Rice coding with parameter \(k = 2\) (i.e., \(m = 4\)):

\[
N = 6
\]

\subsubsection*{Solution}

Compute:

\[
q = N >> k = 6 >> 2 = 1
\]
\[
r = N \& (2^2 - 1) = 6 \& 3 = 6 \bmod 4 = 2
\]

Unary code of \(q=1\):
\[
10
\]

Binary representation of \(r = 2\) with \(k=2\) bits:
\[
10
\]

\subsubsection*{Final Rice Code}

\[
\boxed{10\,10}
\]

\subsection{Why Golomb--Rice Coding Works for Images}

In predictive image coding (JPEG-LS), the residuals:

\[
e(i,j) = x(i,j) - \hat{x}(i,j)
\]

tend to follow a Laplacian distribution centered at zero.  
After mapping to non-negative values using the standard “zig-zag” mapping:

\[
0,1,-1,2,-2,3,-3,\ldots \quad \Rightarrow \quad 0,1,2,3,4,5,6,\ldots
\]

the distribution becomes geometric-like:

\[
P(N = k) \propto (1-p)^k
\]

Golomb and Rice coding are \textbf{optimal} for such distributions.

\subsection{Summary}

\begin{itemize}
\item Golomb codes optimally compress geometric distributions.
\item Code for integer \(N\): unary(\(q\)) + truncated-binary(\(r\)).
\item Rice coding is a special case with \(m=2^k\), making encoding extremely efficient.
\item Widely used in predictive lossless image compression (JPEG-LS).
\item Two steps encode each value: quotient (run-length-like) and remainder (binary-coded).
\end{itemize}

\section{Haar Wavelet Transform}

\subsection{Introduction}

Wavelet transforms provide a multi-resolution representation of signals and images.  
Unlike the Fourier transform, which represents global frequency content, wavelets provide:

\begin{itemize}
\item local frequency information,
\item multi-scale decomposition,
\item sparse representations useful for compression.
\end{itemize}

The \textbf{Haar wavelet} is the simplest orthonormal wavelet, introduced by Alfréd Haar in 1910.  
It is widely used in:

\begin{itemize}
\item image compression (JPEG2000 preliminaries),
\item denoising,
\item feature extraction,
\item multi-resolution analysis.
\end{itemize}

\subsection{Haar Scaling and Wavelet Functions}

The Haar scaling function is:

\[
\phi(t) =
\begin{cases}
1, & 0 \le t < 1, \\
0, & \text{otherwise}.
\end{cases}
\]

The wavelet function is:

\[
\psi(t) =
\begin{cases}
1, & 0 \le t < \frac{1}{2}, \\
-1, & \frac{1}{2} \le t < 1, \\
0, & \text{otherwise}.
\end{cases}
\]

These functions generate the Haar basis for $L^2(\mathbb{R})$.

\subsection{1D Haar Transform}

Given a discrete signal $x = [x_1, x_2, \ldots, x_N]$, with $N$ a power of 2, the 1-level Haar transform is defined as:

\[
A_k = \frac{x_{2k-1} + x_{2k}}{2},
\qquad
D_k = \frac{x_{2k-1} - x_{2k}}{2},
\quad k = 1,\ldots,N/2.
\]

\[
\text{Output} = [A_1, A_2, \ldots, A_{N/2},\; D_1, D_2, \ldots, D_{N/2}]
\]

\subsection{Inverse 1D Haar Transform}

\[
x_{2k-1} = A_k + D_k,
\qquad
x_{2k} = A_k - D_k.
\]

This reconstruction is exact since Haar is orthonormal.

\subsection{2D Haar Transform (Images)}

The 2D Haar transform applies the 1D transform:

\begin{enumerate}
\item to each \textbf{row},
\item then to each \textbf{column}.
\end{enumerate}

The four resulting subbands are:

\[
\begin{bmatrix}
LL & LH \\
HL & HH
\end{bmatrix}
\]

Where:

\begin{itemize}
\item $LL$ = low-pass rows + low-pass columns,
\item $LH$ = low-pass rows + high-pass columns,
\item $HL$ = high-pass rows + low-pass columns,
\item $HH$ = high-pass rows + high-pass columns.
\end{itemize}

\subsection{Wavelet-Based Image Compression}

Compression occurs because:

\begin{itemize}
\item Most energy is in the $LL$ band,
\item Detail bands ($LH$, $HL$, $HH$) often contain small coefficients,
\item Thresholding those coefficients introduces minimal perceptual distortion.
\end{itemize}

Thresholding rule:

\[
c' =
\begin{cases}
c, & |c| \ge T, \\
0, & |c| < T.
\end{cases}
\]

The number of non-zero coefficients measures compression efficiency.

\newpage
\section*{Exam-Level Problems With Full Solutions}

\subsection*{Problem 1: 1D Haar Transform}

Compute the 1-level Haar transform of:
\[
x = [8, 6, 4, 2].
\]

\subsubsection*{Solution}

Approximation coefficients:
\[
A_1 = \frac{8+6}{2} = 7,\qquad A_2 = \frac{4+2}{2} = 3.
\]

Detail coefficients:
\[
D_1 = \frac{8-6}{2} = 1,\qquad D_2 = \frac{4-2}{2} = 1.
\]

\[
\boxed{
[7,\; 3,\; 1,\; 1]
}
\]

---

\subsection*{Problem 2: 2D Haar Transform of Image Block}

Given:
\[
I =
\begin{bmatrix}
4 & 4 & 8 & 8 \\
4 & 4 & 8 & 8 \\
10 & 10 & 6 & 6 \\
10 & 10 & 6 & 6
\end{bmatrix}
\]

Compute 1-level 2D Haar transform.

\subsubsection*{Solution}

\paragraph{Step 1: Row transforms}

Row averages/differences:

Row 1:
\[
[4,4,8,8] \rightarrow [4,8 \;|\; 0,0]
\]

Row 2:
\[
[4,4,8,8] \rightarrow [4,8 \;|\; 0,0]
\]

Row 3:
\[
[10,10,6,6] \rightarrow [10,6 \;|\; 0,0]
\]

Row 4:
\[
[10,10,6,6] \rightarrow [10,6 \;|\; 0,0]
\]

Matrix now:
\[
\begin{bmatrix}
4 & 8 & 0 & 0 \\
4 & 8 & 0 & 0 \\
10 & 6 & 0 & 0 \\
10 & 6 & 0 & 0
\end{bmatrix}
\]

\paragraph{Step 2: Column transforms on left two columns}

Column 1:
\[
[4,4,10,10] \rightarrow [4,10 \;|\; 0,0]
\]

Column 2:
\[
[8,8,6,6] \rightarrow [8,6 \;|\; 0,0]
\]

\paragraph{Final Subbands}

\[
LL =
\begin{bmatrix}
4 & 8 \\
10 & 6
\end{bmatrix},\quad
LH =
\begin{bmatrix}
0 & 0 \\
0 & 0
\end{bmatrix},
\]

\[
HL =
\begin{bmatrix}
0 & 0 \\
0 & 0
\end{bmatrix},\quad
HH =
\begin{bmatrix}
0 & 0 \\
0 & 0
\end{bmatrix}
\]

---

\subsection*{Problem 3: Inverse 1D Haar Transform}

Given:

\[
A_1 = 7,\; A_2 = 3,\qquad
D_1 = 1,\; D_2 = -1
\]

Reconstruct signal.

\subsubsection*{Solution}

\[
x_1 = A_1 + D_1 = 8,\quad x_2 = A_1 - D_1 = 6
\]

\[
x_3 = A_2 + D_2 = 2,\quad x_4 = A_2 - D_2 = 4
\]

\[
\boxed{[8,\; 6,\; 2,\; 4]}
\]

---

\subsection*{Problem 4: Wavelet Compression (Thresholding)}

Given 2D coefficient matrix:

\[
C =
\begin{bmatrix}
18 & 5 & -3 & 1 \\
2 & -1 & 0 & 0 \\
-4 & 1 & 2 & -1 \\
3 & 0 & 0 & 0
\end{bmatrix}
\]

Threshold \(T=2\).

\subsubsection*{Solution}

Apply hard thresholding:

Values with magnitude $<2$ become 0.

\[
C' =
\begin{bmatrix}
18 & 5 & -3 & 0 \\
2 & 0 & 0 & 0 \\
-4 & 0 & 2 & -1 \\
3 & 0 & 0 & 0
\end{bmatrix}
\]

Count nonzero before:
\[
\text{NZ}_{\text{orig}} = 11
\]

Count nonzero after:
\[
\text{NZ}_{\text{new}} = 8
\]

Compression ratio:
\[
\text{CR} = \frac{11}{8} = 1.375
\]

---

\subsection*{Problem 5: 2-Level Haar Transform (1D)}

Signal:
\[
x = [12, 8, 4, 4, 10, 6, 2, 2]
\]

\subsubsection*{Solution}

\paragraph{Level 1}

Pairs:

\[
A_1 = 10,\; D_1 = 2
\]
\[
A_2 = 4,\; D_2 = 0
\]
\[
A_3 = 8,\; D_3 = 2
\]
\[
A_4 = 2,\; D_4 = 0
\]

\[
A_1 = [10,4,8,2],\quad D_1 = [2,0,2,0]
\]

\paragraph{Level 2 on $A_1$}

Pairs:

\[
A_2^{(1)} = 7,\; D_2^{(1)} = 3
\]
\[
A_2^{(2)} = 5,\; D_2^{(2)} = 3
\]

Thus:

\[
A_2 = [7,5],\quad D_2 = [3,3]
\]

\subsection*{Final 2-Level Output}

\[
\boxed{
A_2 = [7,5],\;
D_2 = [3,3],\;
D_1 = [2,0,2,0]
}
\]

\[
\boxed{
[7,\;5,\;3,\;3,\;2,\;0,\;2,\;0]
}
\]

---

\section*{Summary}

\begin{itemize}
\item Haar transform performs averaging and differencing operations.
\item 2D Haar works by applying 1D transform to rows and columns.
\item LL preserves most information, LH/HL/HH capture edges.
\item Thresholding wavelet coefficients yields compression.
\item Multi-level transforms reveal progressively coarser approximations.
\end{itemize}

\section{Wavelet Compression Using the LL Band}

\subsection{Introduction}

Wavelet-based image compression exploits the fact that natural images concentrate most of their visual energy in the low-frequency components.  
After a 2D wavelet transform, the image is decomposed into four subbands:

\[
\begin{bmatrix}
LL & LH \\
HL & HH
\end{bmatrix}
\]

where:

\begin{itemize}
\item $LL$ contains coarse approximation (low frequencies),
\item $LH$ contains horizontal edges,
\item $HL$ contains vertical edges,
\item $HH$ contains diagonal edges.
\end{itemize}

The $LL$ subband retains the majority of the signal energy and is recursively decomposed in multi-level wavelet transforms (as in JPEG2000).

\subsection{Energy Compaction Property}

For many natural signals, wavelet transforms achieve \textbf{sparse representations}.  
Mathematically:

\[
\|LL\|_2^2 \gg \|LH\|_2^2 + \|HL\|_2^2 + \|HH\|_2^2.
\]

That is, most of the energy lies in the approximation band.  
This allows:

\begin{itemize}
\item aggressive quantization of detail bands,
\item preservation of perceptually important low-frequency content,
\item high compression ratios with minimal visual distortion.
\end{itemize}

\subsection{Compression Pipeline}

A typical wavelet-based compression system performs:

\begin{enumerate}
\item \textbf{2D wavelet transform}
\[
I \xrightarrow{W} \{LL, LH, HL, HH\}.
\]

\item \textbf{Quantization}  
For detail coefficients:
\[
c' = \text{sign}(c)\cdot \max(|c| - T,\,0)
\]
or:
\[
c' = 
\begin{cases}
0, & |c| < T, \\
\text{round}(c/Q), & |c| \ge T.
\end{cases}
\]

\item \textbf{Entropy coding} (Huffman, arithmetic, or Golomb–Rice).

\item \textbf{Reconstruction}
\[
\{LL', LH', HL', HH'\} \xrightarrow{W^{-1}} \hat{I}.
\]
\end{enumerate}

\subsection{Multi-Level LL Decomposition}

Recursive decomposition applies the wavelet transform to the $LL$ subband:

\[
LL_0 = I
\]

\[
LL_1, LH_1, HL_1, HH_1 = W(LL_0)
\]

\[
LL_2, LH_2, HL_2, HH_2 = W(LL_1)
\]

After $L$ levels:

\[
\text{Output} = \{LL_L,\; LH_L,\; HL_L,\; HH_L,\ldots, LH_1, HL_1, HH_1\}.
\]

Only $LL_L$ needs to be stored with high fidelity; the rest can be heavily quantized.

\subsection{Rate--Distortion Considerations}

For a given bitrate, optimal bit allocation follows:

\[
b_{LL} > b_{LH} \approx b_{HL} > b_{HH}.
\]

The distortion $D$ is:

\[
D = \sum_{\text{subbands}} w_s \cdot E[(c_s - \hat{c}_s)^2],
\]

where $w_s$ accounts for each band’s contribution to image quality.

\subsection{Compression Ratio Calculation}

If the original image has $N$ coefficients and the compressed representation contains $K$ non-zero or entropy-coded coefficients:

\[
\text{CR} = \frac{N}{K}.
\]

\newpage
\section*{Exam-Level Problems With Complete Solutions}

\subsection*{Problem 1: Energy Compaction in Wavelet Transform}

Given the subband energies:
\[
E_{LL}=8200,\quad E_{LH}=500,\quad E_{HL}=620,\quad E_{HH}=300.
\]

Compute the percentage of total energy in the LL band.

\subsubsection*{Solution}

Total energy:
\[
E_{\text{tot}} = 8200 + 500 + 620 + 300 = 9620.
\]

\[
\%LL = \frac{8200}{9620}\times 100 = 85.2\%.
\]

\[
\boxed{85.2\% \text{ of the energy lies in the LL band.}}
\]

---

\subsection*{Problem 2: Thresholding for Compression}

Given detail coefficients:

\[
D = [3, -1, 0, 2, 7, -3, 1]
\]

Apply hard thresholding with $T=3$.

\subsubsection*{Solution}

Rule:
\[
c' = 
\begin{cases}
0, & |c| < 3 \\
c, & |c| \ge 3
\end{cases}
\]

Apply:

\[
[3, 0, 0, 0, 7, -3, 0]
\]

Number of nonzero coefficients:
\[
K = 3
\]

Original count:
\[
N = 7
\]

Compression ratio:
\[
CR = \frac{7}{3} = 2.33.
\]

---

\subsection*{Problem 3: Multi-Level LL Decomposition}

An image is decomposed using two-level Haar transform:

\[
LL_2 =
\begin{bmatrix}
50 & 60 \\
55 & 65
\end{bmatrix}
\]

\(LH_2, HL_2, HH_2\) have small values and are heavily quantized to zero.

Explain qualitatively what reconstruction looks like.

\subsubsection*{Solution}

Since:

\begin{itemize}
\item $LL_2$ contains coarse approximation of the image,
\item detail coefficients at level 2 are zero,
\item detail coefficients at level 1 are small or zero,
\end{itemize}

Reconstruction approximates a blurred/smoothed image with dominant low-frequency components preserved.

\[
\boxed{\text{Reconstruction looks like a low-resolution, smoothed version of the original image.}}
\]

---

\subsection*{Problem 4: 2D Wavelet Coefficient Storage}

A $256 \times 256$ image undergoes 3-level wavelet transform.

Compute sizes of subbands.

\subsubsection*{Solution}

Level 1:
\[
LL_1: 128\times128,\quad LH_1:128\times128,\; HL_1:128\times128,\; HH_1:128\times128.
\]

Level 2 (applied to $LL_1$):
\[
LL_2: 64\times64,\quad LH_2:64\times64,\; HL_2:64\times64,\; HH_2:64\times64.
\]

Level 3:
\[
LL_3: 32\times32,\quad LH_3:32\times32,\; HL_3:32\times32,\; HH_3:32\times32.
\]

Total coefficient count:
\[
32^2 + 3\cdot 64^2 + 3\cdot 128^2 = 1024 + 12288 + 49152 = 62464,
\]

which equals $256^2 = 65536$ (after rounding behavior differences).

---

\subsection*{Problem 5: Compression Using Only LL Band}

A $512\times512$ image is compressed by storing:

\begin{itemize}
\item the $LL_2$ band (size $128\times128$),
\item 20\% of the largest detail coefficients from all other bands.
\end{itemize}

Compute the fraction of coefficients retained.

\subsubsection*{Solution}

Total coefficients:
\[
= 512^2 = 262144.
\]

LL$_2$ size:
\[
128^2 = 16384.
\]

Detail coefficients total:
\[
262144 - 16384 = 245760.
\]

We keep 20\% of them:
\[
0.20\times245760 = 49152.
\]

Total kept:
\[
16384 + 49152 = 65536.
\]

Fraction retained:
\[
\frac{65536}{262144} = 0.25.
\]

\[
\boxed{25\% \text{ of coefficients retained.}}
\]

---

\section*{Summary}

\begin{itemize}
\item Wavelet compression relies on the dominance of the LL band.
\item Detail bands can be quantized or thresholded with minimal perceptual loss.
\item Multi-level decomposition increases compression efficiency.
\item Energy largely concentrates in low-frequency wavelet components.
\item Compression ratio improves as more detail coefficients are removed.
\end{itemize}

\section{Predictive Coding (DPCM and JPEG-LS / LOCO-I)}

\subsection{Introduction}

Predictive coding exploits statistical redundancy in images by predicting the value of each pixel from its neighbors.  
Rather than encoding raw pixel intensities, we encode the \textbf{prediction error}, which typically has a sharply peaked distribution around zero and is more compressible.

Predictive coding is central to image compression algorithms including:
\begin{itemize}
\item Differential Pulse Code Modulation (DPCM),
\item Lossless JPEG (obsolete),
\item JPEG-LS (LOCO-I) --- the industry standard for lossless image compression,
\item High-quality predictive codecs (e.g., FLIF).
\end{itemize}

The general predictive model is:

\[
\hat{x}(m,n) = P(\mathcal{N})  \quad \text{predictor from neighborhood}
\]

\[
e(m,n) = x(m,n) - \hat{x}(m,n)
\]

Then we encode $e(m,n)$ instead of $x(m,n)$.

During decoding:
\[
\hat{x}(m,n) = \hat{x}(m,n) + e(m,n)
\]

This requires only that the encoder and decoder use the same predictor.

\subsection{DPCM (Differential Pulse Code Modulation)}

DPCM uses linear predictors based on spatial neighbors. Common linear predictors:

\paragraph{1. First-order predictor}
\[
\hat{x}(m,n) = x(m,n-1)
\]

\paragraph{2. Two-dimensional linear predictor}
\[
\hat{x}(m,n) = ax(m,n-1) + bx(m-1,n) + cx(m-1,n-1)
\]

Typical JPEG predictor (lossless mode):

\[
\hat{x}(m,n) = x(m,n-1) + x(m-1,n) - x(m-1,n-1)
\]

\subsection{Prediction Error}

\[
e(m,n) = x(m,n) - \hat{x}(m,n)
\]

If the predictor is good, $e(m,n)$ has:
\begin{itemize}
\item small magnitude,
\item Laplacian or sharply peaked distribution,
\item low entropy.
\end{itemize}

This reduces average code length when entropy coding is applied.

\subsection{Quantization in Predictive Coding}

For \textbf{lossy} predictive coding:

\[
\hat{e}(m,n) = Q(e(m,n))
\]

Decoder reconstructs:

\[
\hat{x}(m,n) = \hat{x}(m,n)+\hat{e}(m,n)
\]

Quantization introduces distortion but reduces entropy dramatically.

\subsection{JPEG-LS (LOCO-I): Nonlinear Predictive Coding}

JPEG-LS uses the \textbf{Median Edge Detector (MED)} predictor.

Let pixels be defined as:

\[
\begin{array}{cc}
C = x(m,n-1), & B = x(m-1,n) \\
A = x(m-1,n-1) &
\end{array}
\]

\begin{center}
\begin{tabular}{ccc}
A & B \\
C & X
\end{tabular}
\end{center}

The MED predictor is:

\[
\hat{x} =
\begin{cases}
\min(B,C), & A \ge \max(B,C), \\
\max(B,C), & A \le \min(B,C), \\
B + C - A, & \text{otherwise}.
\end{cases}
\]

\subsubsection*{Interpretation}

The MED predictor adapts to local gradients:

\begin{itemize}
\item If $A$ is very large (upper-left corner is bright), it predicts the minimum of neighbors: indicating a downward edge.
\item If $A$ is very small, predicts the maximum: indicating an upward edge.
\item Otherwise, uses a linear predictor $B + C - A$ which estimates smooth regions.
\end{itemize}

This nonlinearity allows JPEG-LS to outperform older linear predictors.

\subsection{Run Mode in JPEG-LS}

When the image contains long flat regions:

\[
x(m,n) \approx x(m,n-1) \approx x(m,n-2) \approx \cdots
\]

JPEG-LS switches to run-length mode:

\begin{itemize}
\item Encode the length of the flat run.
\item Encode a terminating sample in regular mode.
\end{itemize}

This exploits the fact that natural images often contain smooth areas.

\subsection{Context Modeling (LOCO-I)}

Prediction residuals are not encoded uniformly.  
Instead, JPEG-LS uses 365 contexts derived from local gradients:

\[
g_1 = B - A,\quad g_2 = A - C,\quad g_3 = C - B
\]

Each context $k$ stores:

\begin{itemize}
\item bias,
\item variance-like estimate,
\item Golomb parameter for entropy coding.
\end{itemize}

This makes JPEG-LS nearly optimal for lossless coding.

\subsection{Entropy Coding in Predictive Coders}

Prediction error $e(m,n)$ is entropy-coded using **Golomb-Rice coding**, ideal for Laplacian distributions.

\[
P(e) \approx \frac{1}{2} \alpha e^{-\alpha|e|}
\]

Small errors occur frequently, so Rice codes are efficient.

\subsection{Rate--Distortion}

For lossy predictive coding:

\[
D = E[(x - \hat{x})^2]
\]

Bitrate increases with:
\[
R \approx H(e_Q)
\]

Optimal quantizer step size balances $R$ and $D$.

\newpage
\section*{Exam-Level Problems With Solutions}

\subsection*{Problem 1: MED Prediction}

Given:
\[
A=120,\; B=100,\; C=90.
\]

Compute the MED prediction.

\subsubsection*{Solution}

Check conditions:

\[
A = 120 \ge \max(B,C) = 100 \Rightarrow \hat{x} = \min(B,C) = 90.
\]

\[
\boxed{\hat{x} = 90}
\]

---

\subsection*{Problem 2: Compute Prediction Error}

Given:
\[
x(m,n)=110,\quad \hat{x}=100.
\]

\[
e = 110 - 100 = 10.
\]

If quantizer $Q(e) = \text{round}(e/4)$:

\[
\hat{e}= 10/4 = 2.5 \rightarrow 3.
\]

Decoder reconstructs:

\[
\hat{x} = 100 + 3 = 103.
\]

\[
\boxed{\text{Reconstructed value } = 103}
\]

---

\subsection*{Problem 3: DPCM Linear Prediction}

Image row:
\[
[100, 104, 108, 130]
\]

Predictor:
\[
\hat{x}(n)=x(n-1)
\]

Compute prediction errors.

\subsubsection*{Solution}

\[
\begin{aligned}
\hat{x}(2) &= 100,\quad e=104-100=4 \\
\hat{x}(3) &= 104,\quad e=108-104=4 \\
\hat{x}(4) &= 108,\quad e=130-108=22
\end{aligned}
\]

\[
\boxed{e = [4,4,22]}
\]

---

\subsection*{Problem 4: JPEG-LS Gradient Context}

Given:
\[
A=50,\; B=60,\; C=55.
\]

Compute the gradients:

\[
g_1 = B-A = 10
\]
\[
g_2 = A-C = -5
\]
\[
g_3 = C-B = -5
\]

This defines the context used for Golomb coding.

\[
\boxed{(g_1,g_2,g_3) = (10,-5,-5)}
\]

---

\subsection*{Problem 5: Compression Gain from Prediction}

Original pixel entropy:
\[
H(X)=7.8\ \text{bits}
\]

Residual entropy:
\[
H(E)=4.1\ \text{bits}
\]

Compute compression gain:

\[
G = \frac{H(X)}{H(E)} = \frac{7.8}{4.1} = 1.90.
\]

\[
\boxed{\text{Prediction reduces coding cost by a factor of } 1.9}
\]

---

\section*{Summary}

\begin{itemize}
\item DPCM uses linear prediction; JPEG-LS uses nonlinear MED prediction.
\item Prediction reduces entropy of image data dramatically.
\item JPEG-LS uses context modeling and Golomb-Rice coding for optimal performance.
\item Run mode efficiently handles flat regions.
\item Predictive coding is essential for both lossless and lossy compression.
\end{itemize}

\section{Morphological Operations}

\subsection{Introduction}

Morphological image processing is based on set theory and is used to analyze and process geometrical structures within an image.  
Morphology is particularly suited for:

\begin{itemize}
\item noise removal,
\item object shape extraction,
\item boundary detection,
\item segmentation preprocessing,
\item skeletonization and thinning.
\end{itemize}

Binary morphology treats images as sets in $\mathbb{Z}^2$, while grayscale morphology generalizes these ideas using max–min operations.

Let a binary image be represented as a set:
\[
A = \{(x,y) : \text{pixel at } (x,y) \text{ is foreground}\}.
\]

Let $B$ be a structuring element (SE), also represented as a set.

\subsection{Translation and Reflection}

\paragraph{Translation:}
\[
B_{(x,y)} = \{(b_x + x,\; b_y + y) : (b_x, b_y) \in B\}.
\]

\paragraph{Reflection:}
\[
\hat{B} = \{(-b_x,-b_y) : (b_x,b_y) \in B\}.
\]

\subsection{Dilation}

Dilation expands objects in a binary image.

\[
A \oplus B = \{ z : (\hat{B})_z \cap A \neq \emptyset \}.
\]

Interpretation:
\begin{itemize}
\item Slides the SE across the image.
\item If any overlap with foreground exists, output pixel becomes 1.
\end{itemize}

\paragraph{Grayscale Dilation}

\[
(f \oplus B)(x,y) = \max_{(s,t) \in B} \{ f(x-s, y-t) + B(s,t) \}.
\]

\subsection{Erosion}

Erosion shrinks objects:

\[
A \ominus B = \{ z : B_z \subseteq A \}.
\]

Interpretation:
\begin{itemize}
\item SE must fit entirely within $A$ for output pixel to be 1.
\end{itemize}

\paragraph{Grayscale Erosion}

\[
(f \ominus B)(x,y) = \min_{(s,t) \in B} \{ f(x+s,y+t) - B(s,t) \}.
\]

\subsection{Opening}

Opening is erosion followed by dilation:

\[
A \circ B = (A \ominus B) \oplus B.
\]

Effects:
\begin{itemize}
\item Removes small objects.
\item Smooths contours.
\end{itemize}

\subsection{Closing}

Closing is dilation followed by erosion:

\[
A \bullet B = (A \oplus B) \ominus B.
\]

Effects:
\begin{itemize}
\item Fills holes.
\item Connects narrow gaps.
\end{itemize}

\subsection{Hit-or-Miss Transform}

Hit-or-miss detects specific shapes within binary images.

Let $B_1$ match the foreground and $B_2$ match background.

\[
A \otimes (B_1, B_2) = (A \ominus B_1) \cap (A^c \ominus B_2).
\]

Used for:
\begin{itemize}
\item Corner detection,
\item Shape recognition,
\item Skeletonization.
\end{itemize}

\subsection{Morphological Gradient}

Measures boundary thickness:

\[
\text{grad}(A) = (A \oplus B) - (A \ominus B).
\]

\paragraph{Internal boundary:}
\[
\partial_i A = A - (A \ominus B)
\]

\paragraph{External boundary:}
\[
\partial_e A = (A \oplus B) - A
\]

\subsection{Top-Hat and Bottom-Hat Transforms}

\paragraph{Top-hat:}
\[
T = A - (A \circ B)
\]
Enhances small bright elements.

\paragraph{Bottom-hat:}
\[
B_h = (A \bullet B) - A
\]
Enhances dark elements.

\subsection{Applications in Image Processing}

\begin{itemize}
\item Removing salt-and-pepper noise (opening).
\item Closing gaps in edges (closing).
\item Extracting boundaries (morphological gradient).
\item Segmentation preprocessing.
\item Shape extraction (hit-or-miss).
\item Feature extraction in license plates, biometrics, documents.
\end{itemize}

\newpage
\section*{Exam-Level Problems With Solutions}

\subsection*{Problem 1: Binary Dilation}

Given binary set:
\[
A = \{(1,1),(1,2),(2,1)\}
\]

Structuring element:
\[
B = \{(0,0),(1,0)\}
\]

Find $A \oplus B$.

\subsubsection*{Solution}

For each $(x,y) \in A$, add SE offsets:

For $(1,1)$: $(1,1),(2,1)$  
For $(1,2)$: $(1,2),(2,2)$  
For $(2,1)$: $(2,1),(3,1)$

Union:
\[
A \oplus B = \{(1,1),(2,1),(1,2),(2,2),(3,1)\}.
\]

---

\subsection*{Problem 2: Erosion Check}

Same sets $A$ and $B$.

Check whether $(2,1)$ belongs to $A \ominus B$.

\subsubsection*{Solution}

Check if $B_{(2,1)}$ fits inside $A$:

\[
B_{(2,1)} = \{(2,1),(3,1)\}.
\]

$(3,1)$ is \textbf{not} in $A$.

Thus:
\[
(2,1) \notin A \ominus B.
\]

---

\subsection*{Problem 3: Opening and Closing}

Compute opening $A \circ B$ and closing $A \bullet B$.

\subsubsection*{Solution}

We already know erosion removes $(1,1)$ and $(1,2)$ because the SE cannot fit.

Erosion result:
\[
A \ominus B = \{(1,1)\}
\]

Dilation on the result:
\[
(1,1) + B = \{(1,1),(2,1)\}
\]

\[
\boxed{A \circ B = \{(1,1),(2,1)\}}
\]

Closing:

Dilation first:
\[
A \oplus B = \{(1,1),(2,1),(1,2),(2,2),(3,1)\}
\]

Erosion:
Only $(1,1)$ and $(2,1)$ remain.

\[
\boxed{A \bullet B = \{(1,1),(2,1)\}}
\]

---

\subsection*{Problem 4: Morphological Gradient}

Let $A$ be a $3\times3$ block:

\[
A =
\begin{bmatrix}
1 &1 &1 \\
1 &1 &1 \\
1 &1 &1
\end{bmatrix}
\]

SE = 3×3 ones.

Dilated $A$ = all ones (no change).

Eroded $A$ = all ones (still fits).

Thus:

\[
\text{grad}(A)=0.
\]

\[
\boxed{\text{A solid block has zero morphological gradient.}}
\]

---

\subsection*{Problem 5: Hit-or-Miss Shape Detection}

Let $A$ contain a corner shape:

\[
A =
\begin{bmatrix}
1 &1 &0 \\
1 &0 &0 \\
0 &0 &0 
\end{bmatrix}
\]

Foreground SE:
\[
B_1=\{(0,0),(1,0),(0,1)\}
\]

Background SE:
\[
B_2=\{(1,1)\}
\]

Check match at the origin.

\subsubsection*{Solution}

Erosion condition:
\[
B_{1} \subseteq A \quad \Rightarrow \text{true.}
\]

Background erosion:
\[
(1,1) \notin A
\]

Thus hit-or-miss = 1.

\[
\boxed{\text{Corner detected.}}
\]

---

\section*{Summary}

\begin{itemize}
\item Dilation expands objects, erosion shrinks them.
\item Opening removes small objects; closing fills holes.
\item Morphological gradient gives object boundaries.
\item Hit-or-miss detects specific shapes.
\item Grayscale morphology generalizes binary operations via min/max.
\end{itemize}




\end{document}

